% ===========================================================================
%  Appendix: Stochastic Processes and Income
%
%  Usage: Place after \appendix in the main document, then
%         % ===========================================================================
%  Appendix: Stochastic Processes and Income
%
%  Usage: Place after \appendix in the main document, then
%         % ===========================================================================
%  Appendix: Stochastic Processes and Income
%
%  Usage: Place after \appendix in the main document, then
%         % ===========================================================================
%  Appendix: Stochastic Processes and Income
%
%  Usage: Place after \appendix in the main document, then
%         \input{appendix_stochastic_processes}
%
%  Assumes booktabs, amsmath, amssymb, graphicx are loaded.
%  Table paths: CareerCosts/tables/publication/stochastic_processes/
%  Figure paths: CareerCosts/figures/stochastic_processes/
% ===========================================================================

% Helper command for figure notes
\providecommand{\figurenotes}[1]{%
  \par\vspace{0.4em}%
  \begin{minipage}{\linewidth}%
    \footnotesize\textit{Notes:}\ #1%
  \end{minipage}%
}


% ===================================================================
\section{Variable Overview}
\label{sec:variable_overview}
% ===================================================================

Table~\ref{tab:variable_overview} summarises the state variables, decision
variables, and continuous states of the life-cycle model.  Agents are women
who enter the model at age~30 and can live up to age~100.  Each period
corresponds to one calendar year.

\begin{table}[htbp]
\centering
\caption{Variable Overview}
\label{tab:variable_overview}
\small
\begin{tabular}{lll}
\toprule
\textbf{Variable Name} & \textbf{Symbol} & \textbf{Possible Values} \\
\midrule
\multicolumn{3}{l}{\textbf{Decisions}} \\
\addlinespace
Labor Supply        & $d_t$           & \{0: Retired, 1: Unemployed, \\
                    &                 & \phantom{\{} 2: Part-Time, 3: Full-Time\} \\
Care Arrangement    & $c_t$           & \{0: No Care, 1: Light Informal Care, \\
                    &                 & \phantom{\{} 2: Intensive Informal Care, \\
                    &                 & \phantom{\{} 3: Combination Care, 4: Formal Care\} \\
Consumption         & $\tilde{c}_t$   & $[0,\, a_t]$ \\
\addlinespace
\midrule
\multicolumn{3}{l}{\textbf{Discrete States}} \\
\addlinespace
Age                 & $t$             & $30,\ldots,100$ \\
Education Type      & $\tau$          & \{Low, High\} \\
Caregiving Type     &                 & \{A: Not Potential Caregiver, \\
                    &                 & \phantom{\{} B: Potential Caregiver\} \\
Already Retired     &                 & \{0, 1\} \\
Partner State       & $p_t$           & \{0: Single, 1: Working-Age Partner, \\
                    &                 & \phantom{\{} 2: Retired Partner\} \\
Health State        & $h_t$           & \{0: Good Health, 1: Bad Health, 2: Dead\} \\
Job Offer           & $o_t$           & \{0: No Offer, 1: Job Offer\} \\
Mother Alive        & $m_t$           & \{0: Alive, 1: Recently Dead, 2: Dead\} \\
Mother ADL          &                 & \{0: No ADL, 1: Light, 2: Intensive\} \\
Care Demand         &                 & \{0: None, 1: Light, 2: Intensive\} \\
\addlinespace
\midrule
\multicolumn{3}{l}{\textbf{Continuous States}} \\
\addlinespace
Assets              & $a_t$           & $\mathbb{R}_{\geq 0}$ \\
Work Experience     & $e_t$           & Projection to interval $[0,\,1]$ \\
Taste Shock         & $\epsilon_t(d_t)$ & GEV i.i.d.\ taste shocks \\
\bottomrule
\end{tabular}

\vspace{0.5em}
\begin{minipage}{\linewidth}
\footnotesize\textit{Notes:}
Key derived variables include pension points $PP(x_t)$, the consumption
equivalence scale $n_t(x_t)$, number of children in the household,
hourly wage $w_t(x_t)$, partner income $y^p_t(x_t)$, total household
income $Y_t(x_t, d_t)$, government benefits $B(x_t, d_t)$, and taxes
$T(x_t, d_t)$---all of which are deterministic functions of the state
variables listed above.
\end{minipage}
\end{table}


% ===================================================================
\section{Auxiliary Markov Processes}
\label{sec:auxiliary_markov}
% ===================================================================

All transition processes in this section are estimated on SOEP-Core panel
data, restricted to the structural estimation sample (2010--2017), unless
noted otherwise.  Estimations are performed separately by education type.


% -------------------------------------------------------------------
\subsection{Partner Transitions}
\label{sec:partner_transitions}
% -------------------------------------------------------------------

The partner state $p_t$ evolves stochastically with transition probabilities
that depend on education, age, and the current partner state---none of which
the agent controls directly.  We impose several restrictions that are
informed by the data and simplify the estimation.  First, agents cannot
transition directly from being single to having a retired partner.  Second,
the probability of becoming single is independent of whether the partner is
currently of working age or retired; that is, separation rates are the same
for both partnered states.  Third, there is no transition from a retired
partner back to a working-age partner, reflecting the absorbing nature of
retirement.  Fourth, all partners are assumed to be retired by age~75;
beyond that age, single agents remain single and partnered agents remain
partnered.  This last assumption proxies for the empirical observation that
most late-life partner-state transitions are attributable to spousal death,
which we capture via the widow pension.

Formally, the transition of the partner state is given by:
\begin{equation}
  \pi(p_{t+1} \mid x_t) \;=\; \Lambda_p\!\bigl(Z_p(x_t)'\,\phi_p\bigr),
  \label{eq:partner_transition}
\end{equation}
where $\Lambda_p$ is the multinomial logistic distribution function,
providing transition probabilities over the three partner states.  The
characteristics in $Z_p(x_t)$ are education, age, and the current partner
state.  We follow a two-step estimation strategy: first, we estimate
marriage and separation probabilities, and second, we estimate the
transition of the partner into retirement.  We estimate partner transitions
separately for each of the two education types.

Table~\ref{tab:partner_transition} reports the estimated transition
probabilities.  Figures~\ref{fig:partner_state_no_partner}--\ref{fig:partner_state_retired}
illustrate the model fit by comparing simulated shares with the empirical
counterparts.

\begin{table}[htbp]
\centering
\caption{Partner Transition Probabilities}
\label{tab:partner_transition}
\input{CareerCosts/tables/publication/stochastic_processes/partner_transition_women}
\end{table}

\begin{figure}[htbp]
\centering
\includegraphics[width=0.85\textwidth]{CareerCosts/figures/stochastic_processes/partner_state_no_partner_women}
\caption{Share of Women Without a Partner}
\label{fig:partner_state_no_partner}
\figurenotes{Observed shares (dashed) and estimated shares (solid) by age
and education type.  Estimation on SOEP-Core, women only.}
\end{figure}

\begin{figure}[htbp]
\centering
\includegraphics[width=0.85\textwidth]{CareerCosts/figures/stochastic_processes/partner_state_working_partner_women}
\caption{Share of Women With a Working-Age Partner}
\label{fig:partner_state_working}
\figurenotes{Observed shares (dashed) and estimated shares (solid) by age
and education type.  Estimation on SOEP-Core, women only.}
\end{figure}

\begin{figure}[htbp]
\centering
\includegraphics[width=0.85\textwidth]{CareerCosts/figures/stochastic_processes/partner_state_retired_partner_women}
\caption{Share of Women With a Retired Partner}
\label{fig:partner_state_retired}
\figurenotes{Observed shares (dashed) and estimated shares (solid) by age
and education type.  Estimation on SOEP-Core, women only.}
\end{figure}


% -------------------------------------------------------------------
\subsection{Number of Children in Household}
\label{sec:children}
% -------------------------------------------------------------------

The agent's partner state, together with age and education type, determines
the number of children living in the household.  The number of children
directly enters the construction of the consumption equivalence scale and,
conditional on the agent being employed, affects utility through the
preference parameters for children.  We approximate the number of children
in the household via OLS, estimating the following specification separately
by education type and partner status:
\begin{equation}
  n_{t}^{\text{child}} \;=\; \alpha_{0} + \alpha_{1}\,\text{period}_{t}
  + \alpha_{2}\,\text{period}_{t}^{2} + u_{t},
  \label{eq:children}
\end{equation}
where $\text{period}_{t} = t - 30$ denotes model age.

Table~\ref{tab:children} presents the OLS estimates.
Figures~\ref{fig:children_single} and~\ref{fig:children_partnered}
show the estimation fit for single and partnered women, respectively.

\begin{table}[htbp]
\centering
\caption{Number of Children in Household --- OLS Parameters}
\label{tab:children}
\input{CareerCosts/tables/publication/stochastic_processes/number_of_children_women}
\end{table}

\begin{figure}[htbp]
\centering
\includegraphics[width=0.85\textwidth]{CareerCosts/figures/stochastic_processes/children_single_women}
\caption{Number of Children --- Single Women}
\label{fig:children_single}
\figurenotes{Observed (dashed) and estimated (solid) mean number of children
in the household by age and education.  SOEP-Core, single women.}
\end{figure}

\begin{figure}[htbp]
\centering
\includegraphics[width=0.85\textwidth]{CareerCosts/figures/stochastic_processes/children_partnered_women}
\caption{Number of Children --- Partnered Women}
\label{fig:children_partnered}
\figurenotes{Observed (dashed) and estimated (solid) mean number of children
in the household by age and education.  SOEP-Core, partnered women.}
\end{figure}


% -------------------------------------------------------------------
\subsection{Job Offer and Separation}
\label{sec:job_offer_separation}
% -------------------------------------------------------------------

The job-offer state $o_t$ governs the agent's ability to choose
employment; the agent can work part-time or full-time only if $o_t = 1$.
If the job-offer state equals zero and the agent has not yet retired, she
is forced into unemployment in the current period.  We incorporate two
independent processes into the job-offer state: job destruction for those
currently employed and job offers for those currently non-employed.

Job destruction probabilities can be directly identified from the SOEP
using the question on the reason for job departure.  We classify
individuals in job-offer state~0 who did not voluntarily leave their
jobs as having experienced an involuntary separation.  The transition
probability for job separation conditional on current employment is given
by:
\begin{equation}
  \pi(o_{t+1} = 0 \mid x_t,\, d_t \in \{2,3\})
  \;=\; \Lambda_{\text{sep}}\!\bigl(Z_{\text{sep}}'\,\phi_{\text{sep}}\bigr),
  \label{eq:job_separation}
\end{equation}
where $\Lambda_{\text{sep}}$ is the logistic distribution function and
$Z_{\text{sep}}$ includes education, age, and a constant.  We separately
estimate separation probabilities by education type.  We restrict the
estimation sample to ages~30--65 to ensure sufficient observational power
and assume that job separation rates remain constant after age~65 until the
mandatory retirement age.

Table~\ref{tab:job_separation} reports the logit estimates and
Figure~\ref{fig:job_separation} shows the predicted separation
probabilities against the observed rates.

\begin{table}[htbp]
\centering
\caption{Job Separation --- Logit Parameters}
\label{tab:job_separation}
\input{CareerCosts/tables/publication/stochastic_processes/job_separation_women}
\end{table}

\begin{figure}[htbp]
\centering
\includegraphics[width=0.85\textwidth]{CareerCosts/figures/stochastic_processes/job_separation_women}
\caption{Job Separation Probabilities}
\label{fig:job_separation}
\figurenotes{Observed (dashed) and estimated (solid) job separation
probabilities by age and education.  Logit estimation using the
\textit{plb0304\_h} variable from the pl sample in the SOEP.}
\end{figure}

The job-offer process for non-employed agents is estimated internally via
the Method of Simulated Moments (MSM).  If the agent chooses non-employment
during a given period, the job-offer process determines the probability of
having an employment choice available in the next period ($o_{t+1} = 1$).


% -------------------------------------------------------------------
\subsection{Health and Death}
\label{sec:health_death}
% -------------------------------------------------------------------

The health process $h_t$ tracks two living states---Good Health and Bad
Health---together with Death.  Survival probabilities are type-specific and
depend on health and age.  This motivates the use of a joint Markov process
to model the stochastic transitions of health and survival simultaneously.

We use the SOEP-Core question on self-reported health to classify
individuals into good and bad health and estimate the probabilities of
transitioning between the two.  Specifically, we use a logistic regression
to estimate the probability of being in bad health in the next period:
\begin{equation}
  \pi(h_{t+1} \mid x_t)
  \;=\; \Lambda_h\!\bigl(Z_h'\,\phi_h\bigr),
  \label{eq:health_transition}
\end{equation}
where $Z_h$ includes the current health state and age.  We estimate the
parameters for each education type separately.  We use the estimated
parameters to predict transition rates and simulate with them from the
initial share of healthy individuals.

Table~\ref{tab:health_transition} documents the estimated parameters.
Figures~\ref{fig:health_good} and~\ref{fig:health_bad} show the predicted
and observed shares of women in good and bad health, respectively.

\begin{table}[htbp]
\centering
\caption{Health Transition Parameters}
\label{tab:health_transition}
\input{CareerCosts/tables/publication/stochastic_processes/health_transition_women}
\end{table}

\begin{figure}[htbp]
\centering
\includegraphics[width=0.85\textwidth]{CareerCosts/figures/stochastic_processes/health_prob_healthy_women}
\caption{Share of Women in Good Health}
\label{fig:health_good}
\figurenotes{Observed (dashed) and estimated (solid) share of women in good
health by age and education.  SOEP-Core.}
\end{figure}

\begin{figure}[htbp]
\centering
\includegraphics[width=0.85\textwidth]{CareerCosts/figures/stochastic_processes/health_prob_bad_women}
\caption{Share of Women in Bad Health}
\label{fig:health_bad}
\figurenotes{Observed (dashed) and estimated (solid) share of women in bad
health by age and education.  SOEP-Core.}
\end{figure}

Death is modeled jointly within the health Markov process $h_t$.  In the
event of death, the agent bequeaths all remaining wealth and receives
bequest utility.  The probability of dying depends on the agent's health
state and education type.  We estimate a logistic specification that scales
federal life-table death probabilities by health--education interaction
terms:
\begin{equation}
  p_{\text{death}}(t, h_t, \tau)
  \;=\; p^{\text{LT}}_{\text{death}}(t)\;\cdot\;\exp\!\bigl(\beta_{h_t,\tau}\bigr),
  \label{eq:mortality}
\end{equation}
where $p^{\text{LT}}_{\text{death}}(t)$ is the baseline death probability
from the Federal Statistical Office life table at age~$t$ and
$\beta_{h_t,\tau}$ is the estimated coefficient for health state $h_t$ and
education type~$\tau$.

Table~\ref{tab:mortality} reports the logit estimates and
Figure~\ref{fig:survival} displays the predicted survival probabilities by
health and education.

\begin{table}[htbp]
\centering
\caption{Mortality --- Logit Parameters}
\label{tab:mortality}
\input{CareerCosts/tables/publication/stochastic_processes/mortality_women}
\end{table}

\begin{figure}[htbp]
\centering
\includegraphics[width=0.85\textwidth]{CareerCosts/figures/stochastic_processes/survival_probabilities_women}
\caption{Share of Women Alive}
\label{fig:survival}
\figurenotes{Estimated survival probabilities by health state and education
type.  Good health (solid) and bad health (dotted).  Life-table baseline
adjusted via logit.  SOEP-Core.}
\end{figure}


% ===================================================================
\section{Parental Health and Care Demand}
\label{sec:parental_health_care}
% ===================================================================

The agent's caregiving decisions depend on her mother's health and care
needs.  We model three interconnected processes: the mother's mortality
(Section~\ref{sec:mother_mortality}), the mother's functional limitations
measured by ADL categories (Section~\ref{sec:mother_adl}), and the
resulting care demand (Section~\ref{sec:care_demand}).  The caregiving type
(A or B) determines whether the agent is a potential informal caregiver;
type-A agents cannot provide informal care themselves, while type-B agents
face the choice between informal and formal care arrangements when care is
demanded.

Additionally, we estimate the probability that another family member
provides informal care when the agent's caregiving type is~A
(Section~\ref{sec:exog_care_supply}).


% -------------------------------------------------------------------
\subsection{Mother's Mortality}
\label{sec:mother_mortality}
% -------------------------------------------------------------------

The mother's survival is modeled via age-specific death probabilities drawn
from federal life tables published by the Federal Statistical Office.  The
mother's age is inferred from the agent's age and the mean
mother--daughter age difference, which we compute by education type from
the SOEP structural sample:
\begin{equation}
  \text{age}^{\text{mother}}_{t}
  \;=\; t + \overline{\Delta}_{\tau}^{\text{mother}} + \delta,
  \label{eq:mother_age}
\end{equation}
where $\overline{\Delta}_{\tau}^{\text{mother}}$ is the mean
mother--daughter age difference for education type~$\tau$ and $\delta$ is a
fixed offset.

The mother's alive/dead state $m_t$ follows a three-state process: Alive
($m_t = 0$), Recently Dead ($m_t = 1$), and Dead ($m_t = 2$).  The
Recently Dead state persists for exactly one period and triggers the
inheritance process (see Section~\ref{sec:inheritance_prob}).  Transitions
are:
\begin{equation}
  P(m_{t+1} \mid m_t) =
  \begin{cases}
    \bigl[1 - p^{\text{LT}}_{\text{death}}(\text{age}^{\text{mother}}_t),\;
           p^{\text{LT}}_{\text{death}}(\text{age}^{\text{mother}}_t),\; 0\bigr]
    & \text{if } m_t = 0, \\[4pt]
    [0,\; 0,\; 1] & \text{if } m_t = 1, \\[4pt]
    [0,\; 0,\; 1] & \text{if } m_t = 2.
  \end{cases}
  \label{eq:mother_death_transition}
\end{equation}

Table~\ref{tab:mother_mortality} reports the life-table parameters used
for the mother's mortality.

\begin{table}[htbp]
\centering
\caption{Mother's Mortality --- Life-Table Parameters}
\label{tab:mother_mortality}
\input{CareerCosts/tables/publication/stochastic_processes/mother_mortality_women}
\end{table}


% -------------------------------------------------------------------
\subsection{Mother's ADL States}
\label{sec:mother_adl}
% -------------------------------------------------------------------

The mother's functional limitations are tracked via Activities of Daily
Living (ADL) categories, which we collapse into three states: No~ADL
(category~0), Light (category~1), and Intensive (categories~2 and~3
combined).  We estimate a multinomial logistic regression on the Survey of
Health, Ageing and Retirement in Europe (SHARE) parent--child data, using
the mother's lagged ADL category and a cubic polynomial in the mother's
age:
\begin{equation}
  P\!\bigl(\text{ADL}_{t+1} = c \mid \text{ADL}_{t},\,\text{age}^{\text{mother}}_t\bigr)
  \;=\;
  \frac{\exp\!\bigl(X_t'\,\beta_c\bigr)}
       {\sum_{j=0}^{C}\exp\!\bigl(X_t'\,\beta_j\bigr)},
  \label{eq:adl_multinomial}
\end{equation}
where $X_t = \bigl(1,\;\text{age},\;\text{age}^2,\;\text{age}^3,\;
\mathbb{1}[\text{ADL}_t = 1],\;\mathbb{1}[\text{ADL}_t = 2],\;
\mathbb{1}[\text{ADL}_t = 3]\bigr)'$
and the reference category is No~ADL ($c = 0$, normalised to
$\beta_0 = 0$).  The model is estimated separately by sex; we use the
female coefficients for the mother.

Table~\ref{tab:mother_adl} reports the multinomial logit estimates.

\begin{table}[htbp]
\centering
\caption{Mother's ADL Transition --- Multinomial Logit Parameters}
\label{tab:mother_adl}
\input{CareerCosts/tables/publication/stochastic_processes/mother_adl_transition_women}
\end{table}


% -------------------------------------------------------------------
\subsection{Care Demand}
\label{sec:care_demand}
% -------------------------------------------------------------------

Care demand is derived deterministically from the mother's ADL state, the
mother's alive/dead status, and a caregiving age window.  Care can only be
demanded while the mother is alive and the agent's age falls within the
caregiving window (ages~40--70).  Outside this window, or when the mother is
dead, care demand is set to zero regardless of the mother's ADL state.

The care demand probabilities are:
\begin{align}
  p_{\text{none}} &= p_{\text{ADL}_0}
    + \bigl(1 - \mathbb{1}_{\text{window}}\bigr)
    \cdot \bigl(p_{\text{ADL}_1} + p_{\text{ADL}_2}\bigr),
  \label{eq:care_demand_none} \\[2pt]
  p_{\text{light}} &= p_{\text{ADL}_1}
    \cdot \mathbb{1}_{\text{window}},
  \label{eq:care_demand_light} \\[2pt]
  p_{\text{intensive}} &= p_{\text{ADL}_2}
    \cdot \mathbb{1}_{\text{window}},
  \label{eq:care_demand_intensive}
\end{align}
where $p_{\text{ADL}_c}$ are the ADL transition probabilities from
equation~\eqref{eq:adl_multinomial} and the indicator
$\mathbb{1}_{\text{window}}$ equals one if and only if the mother is alive
($m_t = 0$) and the agent's age~$t$ lies in $[40,\,70)$.

Since care demand is fully determined by the mother's ADL process and the
caregiving window, no separate estimation is required.


% -------------------------------------------------------------------
\subsection{Exogenous Care Supply}
\label{sec:exog_care_supply}
% -------------------------------------------------------------------

When care is demanded and the agent's caregiving type is~A (not a potential
informal caregiver), another family member may step in to provide informal
care.  We estimate the probability of receiving such informal care from
other family members via a logistic regression on SOEP data:
\begin{equation}
  P(\text{other informal care} = 1)
  \;=\; \Lambda\!\bigl(\gamma_0 + \gamma_1\,\text{age}
  + \gamma_2\,\text{age}^2
  + \gamma_3\,\text{has\_sister}
  + \gamma_4\,\text{education}\bigr),
  \label{eq:exog_care_supply}
\end{equation}
where the sample is restricted to women whose parent has a care demand.

Table~\ref{tab:exog_care_supply} reports the logit estimates.

\begin{table}[htbp]
\centering
\caption{Exogenous Care Supply --- Logit Parameters}
\label{tab:exog_care_supply}
\input{CareerCosts/tables/publication/stochastic_processes/exog_care_supply_women}
\end{table}


% ===================================================================
\section{Income Processes}
\label{sec:income_processes}
% ===================================================================


% -------------------------------------------------------------------
\subsection{Pension Calculation}
\label{sec:pension}
% -------------------------------------------------------------------

The formula for calculating pension claims in Germany consists of three
parts.  First, there is the pension point value, for which we use the
population-weighted average of East and West German pension point values in
2010.  Second, the pension points accumulated over the working life.  Third,
the deduction factor if the individual retired early.  We track the latter
two through the experience stock $e_t$.

Each individual receives pension points in the ratio of her yearly income
compared to the overall mean wage of all working individuals.  Let $w_m$ be
the mean wage and $h_t$ be the agent's work hours (part-time or full-time).
The pension points earned in any given year are:
\begin{equation}
  \text{PP earned}_t \;=\; \frac{w_t \cdot h_t}{w_m}.
  \label{eq:pension_points_earned}
\end{equation}
Accumulated pension points are the sum over all working years.  We
approximate the cumulative pension points by using the mean experience per
type at age~30 and assuming the agent worked full-time thereafter.

If an agent retires at the statutory retirement age (SRA), she receives the
full pension.  The monthly gross pension is:
\begin{equation}
  y_t(x_t, 0) \;=\; PP(x_t) \times PPV,
  \label{eq:pension_gross}
\end{equation}
where $PPV$ is the pension point value (the 2010 population-weighted
East-West average).

If the agent retires before the SRA, an early-retirement penalty of 3.6\%
per year of early retirement applies:
\begin{equation}
  y_t(x_t, 0) \;=\; PP(x_t) \times PPV \times
  \bigl(1 - 0.036 \times \Delta_{\text{age}}\bigr),
  \label{eq:pension_early_retirement}
\end{equation}
where $\Delta_{\text{age}}$ is the number of years between the actual
retirement age and the SRA.  We track the deduction through the experience
stock: the reduced pension corresponds to an adjusted experience level
$e_t'$ such that the early-retirement penalty is embedded in the state
without explicitly tracking the retirement age.


% -------------------------------------------------------------------
\subsection{Own Wage Process}
\label{sec:own_wage}
% -------------------------------------------------------------------

In the model, agents invest in their human capital by working full-time or
part-time.  We estimate the returns to experience using two-way
fixed-effects regressions on SOEP-Core panel data.  The estimation sample
consists of women over age~30 who work full-time or part-time throughout the
estimation period 2010--2017.  Since time fixed effects absorb aggregate
income growth and inflation, all monetary quantities are expressed in 2010
euros.

We estimate the following equation for each education type~$\tau$ using
observations of wages and experience for each individual~$i$ and time~$t$:
\begin{equation}
  \ln(w_{it})
  \;=\; \gamma_{0,\tau} + \gamma_{\ln,\tau}\,\ln(e_{it} + 1)
  + \xi_i + \mu_t + \varepsilon_{it},
  \label{eq:wage}
\end{equation}
where $\xi_i$ denotes individual fixed effects, $\mu_t$ time fixed effects,
and $\varepsilon_{it}$ is the idiosyncratic wage shock.  Our estimates of
$\gamma_{0,\tau}$ and $\gamma_{\ln,\tau}$ directly correspond to the
structural parameters used in the model.  We cluster standard errors across
individuals and time and estimate the wage process's variance
$\sigma_w^2$.

Table~\ref{tab:own_wage} reports the estimates and
Figure~\ref{fig:own_wage} shows the fit of the predicted wages.

\begin{table}[htbp]
\centering
\caption{Own Wage --- Panel OLS Parameters}
\label{tab:own_wage}
\input{CareerCosts/tables/publication/stochastic_processes/own_wage_women}
\end{table}

\begin{figure}[htbp]
\centering
\includegraphics[width=0.85\textwidth]{CareerCosts/figures/stochastic_processes/own_wage_women}
\caption{Own Wage Fit}
\label{fig:own_wage}
\figurenotes{Observed mean log hourly wages (dashed) and predicted values
(solid) by age and education.  Panel OLS with entity and time fixed
effects.  SOEP-Core, women, 2010--2017.}
\end{figure}


% -------------------------------------------------------------------
\subsection{Partner Income Process}
\label{sec:partner_income}
% -------------------------------------------------------------------

We approximate the partner's income through the agent's own state
variables.  If the agent is single ($p_t = 0$), there is no partner income.
When the agent has a working-age partner ($p_t = 1$), we approximate the
partner's monthly gross wage by OLS of wages onto the agent's age and age
squared.  Unemployed partners are assigned a wage of zero.  The partner
income is therefore a combination of the wages of employed partners and the
zero wages of unemployed partners.  We estimate the approximation separately
by education type.

The specification is:
\begin{equation}
  y^p_t \;=\; \alpha_0 + \alpha_1\,\text{age}_t
  + \alpha_2\,\text{age}_t^2 + u_t.
  \label{eq:partner_wage}
\end{equation}

For retired partners ($p_t = 2$), we use the mean income from the
working-age period and apply a replacement rate of 48\% to obtain the
partner's pension income:
\begin{equation}
  y^{p,\text{ret}}_t
  \;=\; 0.48 \times \overline{y^p}.
  \label{eq:partner_pension}
\end{equation}

Table~\ref{tab:partner_wage} presents the OLS estimates and
Figure~\ref{fig:partner_wage} shows the fit of the approximation during the
partner's working life.

\begin{table}[htbp]
\centering
\caption{Partner Wage --- OLS Parameters}
\label{tab:partner_wage}
\input{CareerCosts/tables/publication/stochastic_processes/partner_wage_women}
\end{table}

\begin{figure}[htbp]
\centering
\includegraphics[width=0.85\textwidth]{CareerCosts/figures/stochastic_processes/partner_wage_women}
\caption{Partner Wage Fit}
\label{fig:partner_wage}
\figurenotes{Observed mean monthly gross partner wage (dashed) and OLS
prediction (solid) by age and education.  SOEP-Core, male partners of
women.}
\end{figure}


% ===================================================================
\section{Parental Bequests and Formal Care Costs}
\label{sec:bequests_care_costs}
% ===================================================================


% -------------------------------------------------------------------
\subsection{Inheritance Probability}
\label{sec:inheritance_prob}
% -------------------------------------------------------------------

We model the probability of receiving a positive inheritance conditional
on the mother's recent death ($m_t = 1$).  The inheritance probability
depends on the agent's age, education type, and the care arrangement in
the period of death.  We estimate the following logistic specification:
\begin{equation}
  P(\text{inheritance} > 0)
  \;=\; \Lambda\!\bigl(\beta_0
  + \beta_{\text{age}}\,\text{age}
  + \beta_{\text{age}^2}\,\text{age}^2
  + \beta_{\text{care}}
  + \beta_{\text{edu}}\,\text{education}\bigr),
  \label{eq:inheritance_prob}
\end{equation}
where $\beta_{\text{care}}$ captures the effect of the care arrangement
(no care, formal care, or informal care) on the probability of receiving a
bequest.

Table~\ref{tab:inheritance_prob} reports the logit estimates.

\begin{table}[htbp]
\centering
\caption{Inheritance Probability --- Logit Parameters}
\label{tab:inheritance_prob}
\input{CareerCosts/tables/publication/stochastic_processes/inheritance_probability_women}
\end{table}


% -------------------------------------------------------------------
\subsection{Inheritance Amount}
\label{sec:inheritance_amount}
% -------------------------------------------------------------------

Conditional on receiving a positive inheritance, we model the amount via a
log-linear OLS regression:
\begin{equation}
  \ln(\text{amount})
  \;=\; \beta_0
  + \beta_{\text{age}}\,\text{age}
  + \beta_{\text{age}^2}\,\text{age}^2
  + \beta_{\text{care}}
  + \beta_{\text{edu}}\,\text{education}
  + \varepsilon,
  \label{eq:inheritance_amount}
\end{equation}
where the care-type coefficients distinguish between light informal care,
intensive informal care, and formal care (no care is the reference).

Table~\ref{tab:inheritance_amount} reports the estimates.

\begin{table}[htbp]
\centering
\caption{Inheritance Amount --- OLS Parameters}
\label{tab:inheritance_amount}
\input{CareerCosts/tables/publication/stochastic_processes/inheritance_amount_women}
\end{table}


% -------------------------------------------------------------------
\subsection{Formal Care Costs}
\label{sec:formal_care_costs}
% -------------------------------------------------------------------

When the agent arranges formal care for her parent, she incurs a financial
cost representing her contribution to the parent's care arrangement.  We
estimate the agent's monthly out-of-pocket contribution to formal care
costs using the SOEP Innovation Sample (SOEP-IS), which contains detailed
information on care arrangements and associated expenditures.

We estimate the following OLS specification, pooled across education
types to avoid selection bias---since selection into formal care is likely
correlated with education, conditioning on education would bias the
estimated costs:
\begin{equation}
  \text{costs}_t
  \;=\; \alpha_0 + \alpha_1\,\text{age}_t
  + \alpha_2\,\text{age}_t^2 + \varepsilon_t.
  \label{eq:formal_care_costs}
\end{equation}

Table~\ref{tab:formal_care_costs} reports the OLS estimates and
Figure~\ref{fig:formal_care_costs} shows the predicted costs by age.

\begin{table}[htbp]
\centering
\caption{Formal Care Costs --- OLS Parameters}
\label{tab:formal_care_costs}
\input{CareerCosts/tables/publication/stochastic_processes/formal_care_costs}
\end{table}

\begin{figure}[htbp]
\centering
\includegraphics[width=0.85\textwidth]{CareerCosts/figures/stochastic_processes/formal_care_costs_women}
\caption{Formal Care Costs by Age}
\label{fig:formal_care_costs}
\figurenotes{Observed mean monthly contribution to formal care costs
(dashed) and OLS prediction (solid) by age.  SOEP-IS.  Pooled across
education levels.}
\end{figure}

%
%  Assumes booktabs, amsmath, amssymb, graphicx are loaded.
%  Table paths: CareerCosts/tables/publication/stochastic_processes/
%  Figure paths: CareerCosts/figures/stochastic_processes/
% ===========================================================================

% Helper command for figure notes
\providecommand{\figurenotes}[1]{%
  \par\vspace{0.4em}%
  \begin{minipage}{\linewidth}%
    \footnotesize\textit{Notes:}\ #1%
  \end{minipage}%
}


% ===================================================================
\section{Variable Overview}
\label{sec:variable_overview}
% ===================================================================

Table~\ref{tab:variable_overview} summarises the state variables, decision
variables, and continuous states of the life-cycle model.  Agents are women
who enter the model at age~30 and can live up to age~100.  Each period
corresponds to one calendar year.

\begin{table}[htbp]
\centering
\caption{Variable Overview}
\label{tab:variable_overview}
\small
\begin{tabular}{lll}
\toprule
\textbf{Variable Name} & \textbf{Symbol} & \textbf{Possible Values} \\
\midrule
\multicolumn{3}{l}{\textbf{Decisions}} \\
\addlinespace
Labor Supply        & $d_t$           & \{0: Retired, 1: Unemployed, \\
                    &                 & \phantom{\{} 2: Part-Time, 3: Full-Time\} \\
Care Arrangement    & $c_t$           & \{0: No Care, 1: Light Informal Care, \\
                    &                 & \phantom{\{} 2: Intensive Informal Care, \\
                    &                 & \phantom{\{} 3: Combination Care, 4: Formal Care\} \\
Consumption         & $\tilde{c}_t$   & $[0,\, a_t]$ \\
\addlinespace
\midrule
\multicolumn{3}{l}{\textbf{Discrete States}} \\
\addlinespace
Age                 & $t$             & $30,\ldots,100$ \\
Education Type      & $\tau$          & \{Low, High\} \\
Caregiving Type     &                 & \{A: Not Potential Caregiver, \\
                    &                 & \phantom{\{} B: Potential Caregiver\} \\
Already Retired     &                 & \{0, 1\} \\
Partner State       & $p_t$           & \{0: Single, 1: Working-Age Partner, \\
                    &                 & \phantom{\{} 2: Retired Partner\} \\
Health State        & $h_t$           & \{0: Good Health, 1: Bad Health, 2: Dead\} \\
Job Offer           & $o_t$           & \{0: No Offer, 1: Job Offer\} \\
Mother Alive        & $m_t$           & \{0: Alive, 1: Recently Dead, 2: Dead\} \\
Mother ADL          &                 & \{0: No ADL, 1: Light, 2: Intensive\} \\
Care Demand         &                 & \{0: None, 1: Light, 2: Intensive\} \\
\addlinespace
\midrule
\multicolumn{3}{l}{\textbf{Continuous States}} \\
\addlinespace
Assets              & $a_t$           & $\mathbb{R}_{\geq 0}$ \\
Work Experience     & $e_t$           & Projection to interval $[0,\,1]$ \\
Taste Shock         & $\epsilon_t(d_t)$ & GEV i.i.d.\ taste shocks \\
\bottomrule
\end{tabular}

\vspace{0.5em}
\begin{minipage}{\linewidth}
\footnotesize\textit{Notes:}
Key derived variables include pension points $PP(x_t)$, the consumption
equivalence scale $n_t(x_t)$, number of children in the household,
hourly wage $w_t(x_t)$, partner income $y^p_t(x_t)$, total household
income $Y_t(x_t, d_t)$, government benefits $B(x_t, d_t)$, and taxes
$T(x_t, d_t)$---all of which are deterministic functions of the state
variables listed above.
\end{minipage}
\end{table}


% ===================================================================
\section{Auxiliary Markov Processes}
\label{sec:auxiliary_markov}
% ===================================================================

All transition processes in this section are estimated on SOEP-Core panel
data, restricted to the structural estimation sample (2010--2017), unless
noted otherwise.  Estimations are performed separately by education type.


% -------------------------------------------------------------------
\subsection{Partner Transitions}
\label{sec:partner_transitions}
% -------------------------------------------------------------------

The partner state $p_t$ evolves stochastically with transition probabilities
that depend on education, age, and the current partner state---none of which
the agent controls directly.  We impose several restrictions that are
informed by the data and simplify the estimation.  First, agents cannot
transition directly from being single to having a retired partner.  Second,
the probability of becoming single is independent of whether the partner is
currently of working age or retired; that is, separation rates are the same
for both partnered states.  Third, there is no transition from a retired
partner back to a working-age partner, reflecting the absorbing nature of
retirement.  Fourth, all partners are assumed to be retired by age~75;
beyond that age, single agents remain single and partnered agents remain
partnered.  This last assumption proxies for the empirical observation that
most late-life partner-state transitions are attributable to spousal death,
which we capture via the widow pension.

Formally, the transition of the partner state is given by:
\begin{equation}
  \pi(p_{t+1} \mid x_t) \;=\; \Lambda_p\!\bigl(Z_p(x_t)'\,\phi_p\bigr),
  \label{eq:partner_transition}
\end{equation}
where $\Lambda_p$ is the multinomial logistic distribution function,
providing transition probabilities over the three partner states.  The
characteristics in $Z_p(x_t)$ are education, age, and the current partner
state.  We follow a two-step estimation strategy: first, we estimate
marriage and separation probabilities, and second, we estimate the
transition of the partner into retirement.  We estimate partner transitions
separately for each of the two education types.

Table~\ref{tab:partner_transition} reports the estimated transition
probabilities.  Figures~\ref{fig:partner_state_no_partner}--\ref{fig:partner_state_retired}
illustrate the model fit by comparing simulated shares with the empirical
counterparts.

\begin{table}[htbp]
\centering
\caption{Partner Transition Probabilities}
\label{tab:partner_transition}
\input{CareerCosts/tables/publication/stochastic_processes/partner_transition_women}
\end{table}

\begin{figure}[htbp]
\centering
\includegraphics[width=0.85\textwidth]{CareerCosts/figures/stochastic_processes/partner_state_no_partner_women}
\caption{Share of Women Without a Partner}
\label{fig:partner_state_no_partner}
\figurenotes{Observed shares (dashed) and estimated shares (solid) by age
and education type.  Estimation on SOEP-Core, women only.}
\end{figure}

\begin{figure}[htbp]
\centering
\includegraphics[width=0.85\textwidth]{CareerCosts/figures/stochastic_processes/partner_state_working_partner_women}
\caption{Share of Women With a Working-Age Partner}
\label{fig:partner_state_working}
\figurenotes{Observed shares (dashed) and estimated shares (solid) by age
and education type.  Estimation on SOEP-Core, women only.}
\end{figure}

\begin{figure}[htbp]
\centering
\includegraphics[width=0.85\textwidth]{CareerCosts/figures/stochastic_processes/partner_state_retired_partner_women}
\caption{Share of Women With a Retired Partner}
\label{fig:partner_state_retired}
\figurenotes{Observed shares (dashed) and estimated shares (solid) by age
and education type.  Estimation on SOEP-Core, women only.}
\end{figure}


% -------------------------------------------------------------------
\subsection{Number of Children in Household}
\label{sec:children}
% -------------------------------------------------------------------

The agent's partner state, together with age and education type, determines
the number of children living in the household.  The number of children
directly enters the construction of the consumption equivalence scale and,
conditional on the agent being employed, affects utility through the
preference parameters for children.  We approximate the number of children
in the household via OLS, estimating the following specification separately
by education type and partner status:
\begin{equation}
  n_{t}^{\text{child}} \;=\; \alpha_{0} + \alpha_{1}\,\text{period}_{t}
  + \alpha_{2}\,\text{period}_{t}^{2} + u_{t},
  \label{eq:children}
\end{equation}
where $\text{period}_{t} = t - 30$ denotes model age.

Table~\ref{tab:children} presents the OLS estimates.
Figures~\ref{fig:children_single} and~\ref{fig:children_partnered}
show the estimation fit for single and partnered women, respectively.

\begin{table}[htbp]
\centering
\caption{Number of Children in Household --- OLS Parameters}
\label{tab:children}
\input{CareerCosts/tables/publication/stochastic_processes/number_of_children_women}
\end{table}

\begin{figure}[htbp]
\centering
\includegraphics[width=0.85\textwidth]{CareerCosts/figures/stochastic_processes/children_single_women}
\caption{Number of Children --- Single Women}
\label{fig:children_single}
\figurenotes{Observed (dashed) and estimated (solid) mean number of children
in the household by age and education.  SOEP-Core, single women.}
\end{figure}

\begin{figure}[htbp]
\centering
\includegraphics[width=0.85\textwidth]{CareerCosts/figures/stochastic_processes/children_partnered_women}
\caption{Number of Children --- Partnered Women}
\label{fig:children_partnered}
\figurenotes{Observed (dashed) and estimated (solid) mean number of children
in the household by age and education.  SOEP-Core, partnered women.}
\end{figure}


% -------------------------------------------------------------------
\subsection{Job Offer and Separation}
\label{sec:job_offer_separation}
% -------------------------------------------------------------------

The job-offer state $o_t$ governs the agent's ability to choose
employment; the agent can work part-time or full-time only if $o_t = 1$.
If the job-offer state equals zero and the agent has not yet retired, she
is forced into unemployment in the current period.  We incorporate two
independent processes into the job-offer state: job destruction for those
currently employed and job offers for those currently non-employed.

Job destruction probabilities can be directly identified from the SOEP
using the question on the reason for job departure.  We classify
individuals in job-offer state~0 who did not voluntarily leave their
jobs as having experienced an involuntary separation.  The transition
probability for job separation conditional on current employment is given
by:
\begin{equation}
  \pi(o_{t+1} = 0 \mid x_t,\, d_t \in \{2,3\})
  \;=\; \Lambda_{\text{sep}}\!\bigl(Z_{\text{sep}}'\,\phi_{\text{sep}}\bigr),
  \label{eq:job_separation}
\end{equation}
where $\Lambda_{\text{sep}}$ is the logistic distribution function and
$Z_{\text{sep}}$ includes education, age, and a constant.  We separately
estimate separation probabilities by education type.  We restrict the
estimation sample to ages~30--65 to ensure sufficient observational power
and assume that job separation rates remain constant after age~65 until the
mandatory retirement age.

Table~\ref{tab:job_separation} reports the logit estimates and
Figure~\ref{fig:job_separation} shows the predicted separation
probabilities against the observed rates.

\begin{table}[htbp]
\centering
\caption{Job Separation --- Logit Parameters}
\label{tab:job_separation}
\input{CareerCosts/tables/publication/stochastic_processes/job_separation_women}
\end{table}

\begin{figure}[htbp]
\centering
\includegraphics[width=0.85\textwidth]{CareerCosts/figures/stochastic_processes/job_separation_women}
\caption{Job Separation Probabilities}
\label{fig:job_separation}
\figurenotes{Observed (dashed) and estimated (solid) job separation
probabilities by age and education.  Logit estimation using the
\textit{plb0304\_h} variable from the pl sample in the SOEP.}
\end{figure}

The job-offer process for non-employed agents is estimated internally via
the Method of Simulated Moments (MSM).  If the agent chooses non-employment
during a given period, the job-offer process determines the probability of
having an employment choice available in the next period ($o_{t+1} = 1$).


% -------------------------------------------------------------------
\subsection{Health and Death}
\label{sec:health_death}
% -------------------------------------------------------------------

The health process $h_t$ tracks two living states---Good Health and Bad
Health---together with Death.  Survival probabilities are type-specific and
depend on health and age.  This motivates the use of a joint Markov process
to model the stochastic transitions of health and survival simultaneously.

We use the SOEP-Core question on self-reported health to classify
individuals into good and bad health and estimate the probabilities of
transitioning between the two.  Specifically, we use a logistic regression
to estimate the probability of being in bad health in the next period:
\begin{equation}
  \pi(h_{t+1} \mid x_t)
  \;=\; \Lambda_h\!\bigl(Z_h'\,\phi_h\bigr),
  \label{eq:health_transition}
\end{equation}
where $Z_h$ includes the current health state and age.  We estimate the
parameters for each education type separately.  We use the estimated
parameters to predict transition rates and simulate with them from the
initial share of healthy individuals.

Table~\ref{tab:health_transition} documents the estimated parameters.
Figures~\ref{fig:health_good} and~\ref{fig:health_bad} show the predicted
and observed shares of women in good and bad health, respectively.

\begin{table}[htbp]
\centering
\caption{Health Transition Parameters}
\label{tab:health_transition}
\input{CareerCosts/tables/publication/stochastic_processes/health_transition_women}
\end{table}

\begin{figure}[htbp]
\centering
\includegraphics[width=0.85\textwidth]{CareerCosts/figures/stochastic_processes/health_prob_healthy_women}
\caption{Share of Women in Good Health}
\label{fig:health_good}
\figurenotes{Observed (dashed) and estimated (solid) share of women in good
health by age and education.  SOEP-Core.}
\end{figure}

\begin{figure}[htbp]
\centering
\includegraphics[width=0.85\textwidth]{CareerCosts/figures/stochastic_processes/health_prob_bad_women}
\caption{Share of Women in Bad Health}
\label{fig:health_bad}
\figurenotes{Observed (dashed) and estimated (solid) share of women in bad
health by age and education.  SOEP-Core.}
\end{figure}

Death is modeled jointly within the health Markov process $h_t$.  In the
event of death, the agent bequeaths all remaining wealth and receives
bequest utility.  The probability of dying depends on the agent's health
state and education type.  We estimate a logistic specification that scales
federal life-table death probabilities by health--education interaction
terms:
\begin{equation}
  p_{\text{death}}(t, h_t, \tau)
  \;=\; p^{\text{LT}}_{\text{death}}(t)\;\cdot\;\exp\!\bigl(\beta_{h_t,\tau}\bigr),
  \label{eq:mortality}
\end{equation}
where $p^{\text{LT}}_{\text{death}}(t)$ is the baseline death probability
from the Federal Statistical Office life table at age~$t$ and
$\beta_{h_t,\tau}$ is the estimated coefficient for health state $h_t$ and
education type~$\tau$.

Table~\ref{tab:mortality} reports the logit estimates and
Figure~\ref{fig:survival} displays the predicted survival probabilities by
health and education.

\begin{table}[htbp]
\centering
\caption{Mortality --- Logit Parameters}
\label{tab:mortality}
\input{CareerCosts/tables/publication/stochastic_processes/mortality_women}
\end{table}

\begin{figure}[htbp]
\centering
\includegraphics[width=0.85\textwidth]{CareerCosts/figures/stochastic_processes/survival_probabilities_women}
\caption{Share of Women Alive}
\label{fig:survival}
\figurenotes{Estimated survival probabilities by health state and education
type.  Good health (solid) and bad health (dotted).  Life-table baseline
adjusted via logit.  SOEP-Core.}
\end{figure}


% ===================================================================
\section{Parental Health and Care Demand}
\label{sec:parental_health_care}
% ===================================================================

The agent's caregiving decisions depend on her mother's health and care
needs.  We model three interconnected processes: the mother's mortality
(Section~\ref{sec:mother_mortality}), the mother's functional limitations
measured by ADL categories (Section~\ref{sec:mother_adl}), and the
resulting care demand (Section~\ref{sec:care_demand}).  The caregiving type
(A or B) determines whether the agent is a potential informal caregiver;
type-A agents cannot provide informal care themselves, while type-B agents
face the choice between informal and formal care arrangements when care is
demanded.

Additionally, we estimate the probability that another family member
provides informal care when the agent's caregiving type is~A
(Section~\ref{sec:exog_care_supply}).


% -------------------------------------------------------------------
\subsection{Mother's Mortality}
\label{sec:mother_mortality}
% -------------------------------------------------------------------

The mother's survival is modeled via age-specific death probabilities drawn
from federal life tables published by the Federal Statistical Office.  The
mother's age is inferred from the agent's age and the mean
mother--daughter age difference, which we compute by education type from
the SOEP structural sample:
\begin{equation}
  \text{age}^{\text{mother}}_{t}
  \;=\; t + \overline{\Delta}_{\tau}^{\text{mother}} + \delta,
  \label{eq:mother_age}
\end{equation}
where $\overline{\Delta}_{\tau}^{\text{mother}}$ is the mean
mother--daughter age difference for education type~$\tau$ and $\delta$ is a
fixed offset.

The mother's alive/dead state $m_t$ follows a three-state process: Alive
($m_t = 0$), Recently Dead ($m_t = 1$), and Dead ($m_t = 2$).  The
Recently Dead state persists for exactly one period and triggers the
inheritance process (see Section~\ref{sec:inheritance_prob}).  Transitions
are:
\begin{equation}
  P(m_{t+1} \mid m_t) =
  \begin{cases}
    \bigl[1 - p^{\text{LT}}_{\text{death}}(\text{age}^{\text{mother}}_t),\;
           p^{\text{LT}}_{\text{death}}(\text{age}^{\text{mother}}_t),\; 0\bigr]
    & \text{if } m_t = 0, \\[4pt]
    [0,\; 0,\; 1] & \text{if } m_t = 1, \\[4pt]
    [0,\; 0,\; 1] & \text{if } m_t = 2.
  \end{cases}
  \label{eq:mother_death_transition}
\end{equation}

Table~\ref{tab:mother_mortality} reports the life-table parameters used
for the mother's mortality.

\begin{table}[htbp]
\centering
\caption{Mother's Mortality --- Life-Table Parameters}
\label{tab:mother_mortality}
\input{CareerCosts/tables/publication/stochastic_processes/mother_mortality_women}
\end{table}


% -------------------------------------------------------------------
\subsection{Mother's ADL States}
\label{sec:mother_adl}
% -------------------------------------------------------------------

The mother's functional limitations are tracked via Activities of Daily
Living (ADL) categories, which we collapse into three states: No~ADL
(category~0), Light (category~1), and Intensive (categories~2 and~3
combined).  We estimate a multinomial logistic regression on the Survey of
Health, Ageing and Retirement in Europe (SHARE) parent--child data, using
the mother's lagged ADL category and a cubic polynomial in the mother's
age:
\begin{equation}
  P\!\bigl(\text{ADL}_{t+1} = c \mid \text{ADL}_{t},\,\text{age}^{\text{mother}}_t\bigr)
  \;=\;
  \frac{\exp\!\bigl(X_t'\,\beta_c\bigr)}
       {\sum_{j=0}^{C}\exp\!\bigl(X_t'\,\beta_j\bigr)},
  \label{eq:adl_multinomial}
\end{equation}
where $X_t = \bigl(1,\;\text{age},\;\text{age}^2,\;\text{age}^3,\;
\mathbb{1}[\text{ADL}_t = 1],\;\mathbb{1}[\text{ADL}_t = 2],\;
\mathbb{1}[\text{ADL}_t = 3]\bigr)'$
and the reference category is No~ADL ($c = 0$, normalised to
$\beta_0 = 0$).  The model is estimated separately by sex; we use the
female coefficients for the mother.

Table~\ref{tab:mother_adl} reports the multinomial logit estimates.

\begin{table}[htbp]
\centering
\caption{Mother's ADL Transition --- Multinomial Logit Parameters}
\label{tab:mother_adl}
\input{CareerCosts/tables/publication/stochastic_processes/mother_adl_transition_women}
\end{table}


% -------------------------------------------------------------------
\subsection{Care Demand}
\label{sec:care_demand}
% -------------------------------------------------------------------

Care demand is derived deterministically from the mother's ADL state, the
mother's alive/dead status, and a caregiving age window.  Care can only be
demanded while the mother is alive and the agent's age falls within the
caregiving window (ages~40--70).  Outside this window, or when the mother is
dead, care demand is set to zero regardless of the mother's ADL state.

The care demand probabilities are:
\begin{align}
  p_{\text{none}} &= p_{\text{ADL}_0}
    + \bigl(1 - \mathbb{1}_{\text{window}}\bigr)
    \cdot \bigl(p_{\text{ADL}_1} + p_{\text{ADL}_2}\bigr),
  \label{eq:care_demand_none} \\[2pt]
  p_{\text{light}} &= p_{\text{ADL}_1}
    \cdot \mathbb{1}_{\text{window}},
  \label{eq:care_demand_light} \\[2pt]
  p_{\text{intensive}} &= p_{\text{ADL}_2}
    \cdot \mathbb{1}_{\text{window}},
  \label{eq:care_demand_intensive}
\end{align}
where $p_{\text{ADL}_c}$ are the ADL transition probabilities from
equation~\eqref{eq:adl_multinomial} and the indicator
$\mathbb{1}_{\text{window}}$ equals one if and only if the mother is alive
($m_t = 0$) and the agent's age~$t$ lies in $[40,\,70)$.

Since care demand is fully determined by the mother's ADL process and the
caregiving window, no separate estimation is required.


% -------------------------------------------------------------------
\subsection{Exogenous Care Supply}
\label{sec:exog_care_supply}
% -------------------------------------------------------------------

When care is demanded and the agent's caregiving type is~A (not a potential
informal caregiver), another family member may step in to provide informal
care.  We estimate the probability of receiving such informal care from
other family members via a logistic regression on SOEP data:
\begin{equation}
  P(\text{other informal care} = 1)
  \;=\; \Lambda\!\bigl(\gamma_0 + \gamma_1\,\text{age}
  + \gamma_2\,\text{age}^2
  + \gamma_3\,\text{has\_sister}
  + \gamma_4\,\text{education}\bigr),
  \label{eq:exog_care_supply}
\end{equation}
where the sample is restricted to women whose parent has a care demand.

Table~\ref{tab:exog_care_supply} reports the logit estimates.

\begin{table}[htbp]
\centering
\caption{Exogenous Care Supply --- Logit Parameters}
\label{tab:exog_care_supply}
\input{CareerCosts/tables/publication/stochastic_processes/exog_care_supply_women}
\end{table}


% ===================================================================
\section{Income Processes}
\label{sec:income_processes}
% ===================================================================


% -------------------------------------------------------------------
\subsection{Pension Calculation}
\label{sec:pension}
% -------------------------------------------------------------------

The formula for calculating pension claims in Germany consists of three
parts.  First, there is the pension point value, for which we use the
population-weighted average of East and West German pension point values in
2010.  Second, the pension points accumulated over the working life.  Third,
the deduction factor if the individual retired early.  We track the latter
two through the experience stock $e_t$.

Each individual receives pension points in the ratio of her yearly income
compared to the overall mean wage of all working individuals.  Let $w_m$ be
the mean wage and $h_t$ be the agent's work hours (part-time or full-time).
The pension points earned in any given year are:
\begin{equation}
  \text{PP earned}_t \;=\; \frac{w_t \cdot h_t}{w_m}.
  \label{eq:pension_points_earned}
\end{equation}
Accumulated pension points are the sum over all working years.  We
approximate the cumulative pension points by using the mean experience per
type at age~30 and assuming the agent worked full-time thereafter.

If an agent retires at the statutory retirement age (SRA), she receives the
full pension.  The monthly gross pension is:
\begin{equation}
  y_t(x_t, 0) \;=\; PP(x_t) \times PPV,
  \label{eq:pension_gross}
\end{equation}
where $PPV$ is the pension point value (the 2010 population-weighted
East-West average).

If the agent retires before the SRA, an early-retirement penalty of 3.6\%
per year of early retirement applies:
\begin{equation}
  y_t(x_t, 0) \;=\; PP(x_t) \times PPV \times
  \bigl(1 - 0.036 \times \Delta_{\text{age}}\bigr),
  \label{eq:pension_early_retirement}
\end{equation}
where $\Delta_{\text{age}}$ is the number of years between the actual
retirement age and the SRA.  We track the deduction through the experience
stock: the reduced pension corresponds to an adjusted experience level
$e_t'$ such that the early-retirement penalty is embedded in the state
without explicitly tracking the retirement age.


% -------------------------------------------------------------------
\subsection{Own Wage Process}
\label{sec:own_wage}
% -------------------------------------------------------------------

In the model, agents invest in their human capital by working full-time or
part-time.  We estimate the returns to experience using two-way
fixed-effects regressions on SOEP-Core panel data.  The estimation sample
consists of women over age~30 who work full-time or part-time throughout the
estimation period 2010--2017.  Since time fixed effects absorb aggregate
income growth and inflation, all monetary quantities are expressed in 2010
euros.

We estimate the following equation for each education type~$\tau$ using
observations of wages and experience for each individual~$i$ and time~$t$:
\begin{equation}
  \ln(w_{it})
  \;=\; \gamma_{0,\tau} + \gamma_{\ln,\tau}\,\ln(e_{it} + 1)
  + \xi_i + \mu_t + \varepsilon_{it},
  \label{eq:wage}
\end{equation}
where $\xi_i$ denotes individual fixed effects, $\mu_t$ time fixed effects,
and $\varepsilon_{it}$ is the idiosyncratic wage shock.  Our estimates of
$\gamma_{0,\tau}$ and $\gamma_{\ln,\tau}$ directly correspond to the
structural parameters used in the model.  We cluster standard errors across
individuals and time and estimate the wage process's variance
$\sigma_w^2$.

Table~\ref{tab:own_wage} reports the estimates and
Figure~\ref{fig:own_wage} shows the fit of the predicted wages.

\begin{table}[htbp]
\centering
\caption{Own Wage --- Panel OLS Parameters}
\label{tab:own_wage}
\input{CareerCosts/tables/publication/stochastic_processes/own_wage_women}
\end{table}

\begin{figure}[htbp]
\centering
\includegraphics[width=0.85\textwidth]{CareerCosts/figures/stochastic_processes/own_wage_women}
\caption{Own Wage Fit}
\label{fig:own_wage}
\figurenotes{Observed mean log hourly wages (dashed) and predicted values
(solid) by age and education.  Panel OLS with entity and time fixed
effects.  SOEP-Core, women, 2010--2017.}
\end{figure}


% -------------------------------------------------------------------
\subsection{Partner Income Process}
\label{sec:partner_income}
% -------------------------------------------------------------------

We approximate the partner's income through the agent's own state
variables.  If the agent is single ($p_t = 0$), there is no partner income.
When the agent has a working-age partner ($p_t = 1$), we approximate the
partner's monthly gross wage by OLS of wages onto the agent's age and age
squared.  Unemployed partners are assigned a wage of zero.  The partner
income is therefore a combination of the wages of employed partners and the
zero wages of unemployed partners.  We estimate the approximation separately
by education type.

The specification is:
\begin{equation}
  y^p_t \;=\; \alpha_0 + \alpha_1\,\text{age}_t
  + \alpha_2\,\text{age}_t^2 + u_t.
  \label{eq:partner_wage}
\end{equation}

For retired partners ($p_t = 2$), we use the mean income from the
working-age period and apply a replacement rate of 48\% to obtain the
partner's pension income:
\begin{equation}
  y^{p,\text{ret}}_t
  \;=\; 0.48 \times \overline{y^p}.
  \label{eq:partner_pension}
\end{equation}

Table~\ref{tab:partner_wage} presents the OLS estimates and
Figure~\ref{fig:partner_wage} shows the fit of the approximation during the
partner's working life.

\begin{table}[htbp]
\centering
\caption{Partner Wage --- OLS Parameters}
\label{tab:partner_wage}
\input{CareerCosts/tables/publication/stochastic_processes/partner_wage_women}
\end{table}

\begin{figure}[htbp]
\centering
\includegraphics[width=0.85\textwidth]{CareerCosts/figures/stochastic_processes/partner_wage_women}
\caption{Partner Wage Fit}
\label{fig:partner_wage}
\figurenotes{Observed mean monthly gross partner wage (dashed) and OLS
prediction (solid) by age and education.  SOEP-Core, male partners of
women.}
\end{figure}


% ===================================================================
\section{Parental Bequests and Formal Care Costs}
\label{sec:bequests_care_costs}
% ===================================================================


% -------------------------------------------------------------------
\subsection{Inheritance Probability}
\label{sec:inheritance_prob}
% -------------------------------------------------------------------

We model the probability of receiving a positive inheritance conditional
on the mother's recent death ($m_t = 1$).  The inheritance probability
depends on the agent's age, education type, and the care arrangement in
the period of death.  We estimate the following logistic specification:
\begin{equation}
  P(\text{inheritance} > 0)
  \;=\; \Lambda\!\bigl(\beta_0
  + \beta_{\text{age}}\,\text{age}
  + \beta_{\text{age}^2}\,\text{age}^2
  + \beta_{\text{care}}
  + \beta_{\text{edu}}\,\text{education}\bigr),
  \label{eq:inheritance_prob}
\end{equation}
where $\beta_{\text{care}}$ captures the effect of the care arrangement
(no care, formal care, or informal care) on the probability of receiving a
bequest.

Table~\ref{tab:inheritance_prob} reports the logit estimates.

\begin{table}[htbp]
\centering
\caption{Inheritance Probability --- Logit Parameters}
\label{tab:inheritance_prob}
\input{CareerCosts/tables/publication/stochastic_processes/inheritance_probability_women}
\end{table}


% -------------------------------------------------------------------
\subsection{Inheritance Amount}
\label{sec:inheritance_amount}
% -------------------------------------------------------------------

Conditional on receiving a positive inheritance, we model the amount via a
log-linear OLS regression:
\begin{equation}
  \ln(\text{amount})
  \;=\; \beta_0
  + \beta_{\text{age}}\,\text{age}
  + \beta_{\text{age}^2}\,\text{age}^2
  + \beta_{\text{care}}
  + \beta_{\text{edu}}\,\text{education}
  + \varepsilon,
  \label{eq:inheritance_amount}
\end{equation}
where the care-type coefficients distinguish between light informal care,
intensive informal care, and formal care (no care is the reference).

Table~\ref{tab:inheritance_amount} reports the estimates.

\begin{table}[htbp]
\centering
\caption{Inheritance Amount --- OLS Parameters}
\label{tab:inheritance_amount}
\input{CareerCosts/tables/publication/stochastic_processes/inheritance_amount_women}
\end{table}


% -------------------------------------------------------------------
\subsection{Formal Care Costs}
\label{sec:formal_care_costs}
% -------------------------------------------------------------------

When the agent arranges formal care for her parent, she incurs a financial
cost representing her contribution to the parent's care arrangement.  We
estimate the agent's monthly out-of-pocket contribution to formal care
costs using the SOEP Innovation Sample (SOEP-IS), which contains detailed
information on care arrangements and associated expenditures.

We estimate the following OLS specification, pooled across education
types to avoid selection bias---since selection into formal care is likely
correlated with education, conditioning on education would bias the
estimated costs:
\begin{equation}
  \text{costs}_t
  \;=\; \alpha_0 + \alpha_1\,\text{age}_t
  + \alpha_2\,\text{age}_t^2 + \varepsilon_t.
  \label{eq:formal_care_costs}
\end{equation}

Table~\ref{tab:formal_care_costs} reports the OLS estimates and
Figure~\ref{fig:formal_care_costs} shows the predicted costs by age.

\begin{table}[htbp]
\centering
\caption{Formal Care Costs --- OLS Parameters}
\label{tab:formal_care_costs}
\input{CareerCosts/tables/publication/stochastic_processes/formal_care_costs}
\end{table}

\begin{figure}[htbp]
\centering
\includegraphics[width=0.85\textwidth]{CareerCosts/figures/stochastic_processes/formal_care_costs_women}
\caption{Formal Care Costs by Age}
\label{fig:formal_care_costs}
\figurenotes{Observed mean monthly contribution to formal care costs
(dashed) and OLS prediction (solid) by age.  SOEP-IS.  Pooled across
education levels.}
\end{figure}

%
%  Assumes booktabs, amsmath, amssymb, graphicx are loaded.
%  Table paths: CareerCosts/tables/publication/stochastic_processes/
%  Figure paths: CareerCosts/figures/stochastic_processes/
% ===========================================================================

% Helper command for figure notes
\providecommand{\figurenotes}[1]{%
  \par\vspace{0.4em}%
  \begin{minipage}{\linewidth}%
    \footnotesize\textit{Notes:}\ #1%
  \end{minipage}%
}


% ===================================================================
\section{Variable Overview}
\label{sec:variable_overview}
% ===================================================================

Table~\ref{tab:variable_overview} summarises the state variables, decision
variables, and continuous states of the life-cycle model.  Agents are women
who enter the model at age~30 and can live up to age~100.  Each period
corresponds to one calendar year.

\begin{table}[htbp]
\centering
\caption{Variable Overview}
\label{tab:variable_overview}
\small
\begin{tabular}{lll}
\toprule
\textbf{Variable Name} & \textbf{Symbol} & \textbf{Possible Values} \\
\midrule
\multicolumn{3}{l}{\textbf{Decisions}} \\
\addlinespace
Labor Supply        & $d_t$           & \{0: Retired, 1: Unemployed, \\
                    &                 & \phantom{\{} 2: Part-Time, 3: Full-Time\} \\
Care Arrangement    & $c_t$           & \{0: No Care, 1: Light Informal Care, \\
                    &                 & \phantom{\{} 2: Intensive Informal Care, \\
                    &                 & \phantom{\{} 3: Combination Care, 4: Formal Care\} \\
Consumption         & $\tilde{c}_t$   & $[0,\, a_t]$ \\
\addlinespace
\midrule
\multicolumn{3}{l}{\textbf{Discrete States}} \\
\addlinespace
Age                 & $t$             & $30,\ldots,100$ \\
Education Type      & $\tau$          & \{Low, High\} \\
Caregiving Type     &                 & \{A: Not Potential Caregiver, \\
                    &                 & \phantom{\{} B: Potential Caregiver\} \\
Already Retired     &                 & \{0, 1\} \\
Partner State       & $p_t$           & \{0: Single, 1: Working-Age Partner, \\
                    &                 & \phantom{\{} 2: Retired Partner\} \\
Health State        & $h_t$           & \{0: Good Health, 1: Bad Health, 2: Dead\} \\
Job Offer           & $o_t$           & \{0: No Offer, 1: Job Offer\} \\
Mother Alive        & $m_t$           & \{0: Alive, 1: Recently Dead, 2: Dead\} \\
Mother ADL          &                 & \{0: No ADL, 1: Light, 2: Intensive\} \\
Care Demand         &                 & \{0: None, 1: Light, 2: Intensive\} \\
\addlinespace
\midrule
\multicolumn{3}{l}{\textbf{Continuous States}} \\
\addlinespace
Assets              & $a_t$           & $\mathbb{R}_{\geq 0}$ \\
Work Experience     & $e_t$           & Projection to interval $[0,\,1]$ \\
Taste Shock         & $\epsilon_t(d_t)$ & GEV i.i.d.\ taste shocks \\
\bottomrule
\end{tabular}

\vspace{0.5em}
\begin{minipage}{\linewidth}
\footnotesize\textit{Notes:}
Key derived variables include pension points $PP(x_t)$, the consumption
equivalence scale $n_t(x_t)$, number of children in the household,
hourly wage $w_t(x_t)$, partner income $y^p_t(x_t)$, total household
income $Y_t(x_t, d_t)$, government benefits $B(x_t, d_t)$, and taxes
$T(x_t, d_t)$---all of which are deterministic functions of the state
variables listed above.
\end{minipage}
\end{table}


% ===================================================================
\section{Auxiliary Markov Processes}
\label{sec:auxiliary_markov}
% ===================================================================

All transition processes in this section are estimated on SOEP-Core panel
data, restricted to the structural estimation sample (2010--2017), unless
noted otherwise.  Estimations are performed separately by education type.


% -------------------------------------------------------------------
\subsection{Partner Transitions}
\label{sec:partner_transitions}
% -------------------------------------------------------------------

The partner state $p_t$ evolves stochastically with transition probabilities
that depend on education, age, and the current partner state---none of which
the agent controls directly.  We impose several restrictions that are
informed by the data and simplify the estimation.  First, agents cannot
transition directly from being single to having a retired partner.  Second,
the probability of becoming single is independent of whether the partner is
currently of working age or retired; that is, separation rates are the same
for both partnered states.  Third, there is no transition from a retired
partner back to a working-age partner, reflecting the absorbing nature of
retirement.  Fourth, all partners are assumed to be retired by age~75;
beyond that age, single agents remain single and partnered agents remain
partnered.  This last assumption proxies for the empirical observation that
most late-life partner-state transitions are attributable to spousal death,
which we capture via the widow pension.

Formally, the transition of the partner state is given by:
\begin{equation}
  \pi(p_{t+1} \mid x_t) \;=\; \Lambda_p\!\bigl(Z_p(x_t)'\,\phi_p\bigr),
  \label{eq:partner_transition}
\end{equation}
where $\Lambda_p$ is the multinomial logistic distribution function,
providing transition probabilities over the three partner states.  The
characteristics in $Z_p(x_t)$ are education, age, and the current partner
state.  We follow a two-step estimation strategy: first, we estimate
marriage and separation probabilities, and second, we estimate the
transition of the partner into retirement.  We estimate partner transitions
separately for each of the two education types.

Table~\ref{tab:partner_transition} reports the estimated transition
probabilities.  Figures~\ref{fig:partner_state_no_partner}--\ref{fig:partner_state_retired}
illustrate the model fit by comparing simulated shares with the empirical
counterparts.

\begin{table}[htbp]
\centering
\caption{Partner Transition Probabilities}
\label{tab:partner_transition}
\input{CareerCosts/tables/publication/stochastic_processes/partner_transition_women}
\end{table}

\begin{figure}[htbp]
\centering
\includegraphics[width=0.85\textwidth]{CareerCosts/figures/stochastic_processes/partner_state_no_partner_women}
\caption{Share of Women Without a Partner}
\label{fig:partner_state_no_partner}
\figurenotes{Observed shares (dashed) and estimated shares (solid) by age
and education type.  Estimation on SOEP-Core, women only.}
\end{figure}

\begin{figure}[htbp]
\centering
\includegraphics[width=0.85\textwidth]{CareerCosts/figures/stochastic_processes/partner_state_working_partner_women}
\caption{Share of Women With a Working-Age Partner}
\label{fig:partner_state_working}
\figurenotes{Observed shares (dashed) and estimated shares (solid) by age
and education type.  Estimation on SOEP-Core, women only.}
\end{figure}

\begin{figure}[htbp]
\centering
\includegraphics[width=0.85\textwidth]{CareerCosts/figures/stochastic_processes/partner_state_retired_partner_women}
\caption{Share of Women With a Retired Partner}
\label{fig:partner_state_retired}
\figurenotes{Observed shares (dashed) and estimated shares (solid) by age
and education type.  Estimation on SOEP-Core, women only.}
\end{figure}


% -------------------------------------------------------------------
\subsection{Number of Children in Household}
\label{sec:children}
% -------------------------------------------------------------------

The agent's partner state, together with age and education type, determines
the number of children living in the household.  The number of children
directly enters the construction of the consumption equivalence scale and,
conditional on the agent being employed, affects utility through the
preference parameters for children.  We approximate the number of children
in the household via OLS, estimating the following specification separately
by education type and partner status:
\begin{equation}
  n_{t}^{\text{child}} \;=\; \alpha_{0} + \alpha_{1}\,\text{period}_{t}
  + \alpha_{2}\,\text{period}_{t}^{2} + u_{t},
  \label{eq:children}
\end{equation}
where $\text{period}_{t} = t - 30$ denotes model age.

Table~\ref{tab:children} presents the OLS estimates.
Figures~\ref{fig:children_single} and~\ref{fig:children_partnered}
show the estimation fit for single and partnered women, respectively.

\begin{table}[htbp]
\centering
\caption{Number of Children in Household --- OLS Parameters}
\label{tab:children}
\input{CareerCosts/tables/publication/stochastic_processes/number_of_children_women}
\end{table}

\begin{figure}[htbp]
\centering
\includegraphics[width=0.85\textwidth]{CareerCosts/figures/stochastic_processes/children_single_women}
\caption{Number of Children --- Single Women}
\label{fig:children_single}
\figurenotes{Observed (dashed) and estimated (solid) mean number of children
in the household by age and education.  SOEP-Core, single women.}
\end{figure}

\begin{figure}[htbp]
\centering
\includegraphics[width=0.85\textwidth]{CareerCosts/figures/stochastic_processes/children_partnered_women}
\caption{Number of Children --- Partnered Women}
\label{fig:children_partnered}
\figurenotes{Observed (dashed) and estimated (solid) mean number of children
in the household by age and education.  SOEP-Core, partnered women.}
\end{figure}


% -------------------------------------------------------------------
\subsection{Job Offer and Separation}
\label{sec:job_offer_separation}
% -------------------------------------------------------------------

The job-offer state $o_t$ governs the agent's ability to choose
employment; the agent can work part-time or full-time only if $o_t = 1$.
If the job-offer state equals zero and the agent has not yet retired, she
is forced into unemployment in the current period.  We incorporate two
independent processes into the job-offer state: job destruction for those
currently employed and job offers for those currently non-employed.

Job destruction probabilities can be directly identified from the SOEP
using the question on the reason for job departure.  We classify
individuals in job-offer state~0 who did not voluntarily leave their
jobs as having experienced an involuntary separation.  The transition
probability for job separation conditional on current employment is given
by:
\begin{equation}
  \pi(o_{t+1} = 0 \mid x_t,\, d_t \in \{2,3\})
  \;=\; \Lambda_{\text{sep}}\!\bigl(Z_{\text{sep}}'\,\phi_{\text{sep}}\bigr),
  \label{eq:job_separation}
\end{equation}
where $\Lambda_{\text{sep}}$ is the logistic distribution function and
$Z_{\text{sep}}$ includes education, age, and a constant.  We separately
estimate separation probabilities by education type.  We restrict the
estimation sample to ages~30--65 to ensure sufficient observational power
and assume that job separation rates remain constant after age~65 until the
mandatory retirement age.

Table~\ref{tab:job_separation} reports the logit estimates and
Figure~\ref{fig:job_separation} shows the predicted separation
probabilities against the observed rates.

\begin{table}[htbp]
\centering
\caption{Job Separation --- Logit Parameters}
\label{tab:job_separation}
\input{CareerCosts/tables/publication/stochastic_processes/job_separation_women}
\end{table}

\begin{figure}[htbp]
\centering
\includegraphics[width=0.85\textwidth]{CareerCosts/figures/stochastic_processes/job_separation_women}
\caption{Job Separation Probabilities}
\label{fig:job_separation}
\figurenotes{Observed (dashed) and estimated (solid) job separation
probabilities by age and education.  Logit estimation using the
\textit{plb0304\_h} variable from the pl sample in the SOEP.}
\end{figure}

The job-offer process for non-employed agents is estimated internally via
the Method of Simulated Moments (MSM).  If the agent chooses non-employment
during a given period, the job-offer process determines the probability of
having an employment choice available in the next period ($o_{t+1} = 1$).


% -------------------------------------------------------------------
\subsection{Health and Death}
\label{sec:health_death}
% -------------------------------------------------------------------

The health process $h_t$ tracks two living states---Good Health and Bad
Health---together with Death.  Survival probabilities are type-specific and
depend on health and age.  This motivates the use of a joint Markov process
to model the stochastic transitions of health and survival simultaneously.

We use the SOEP-Core question on self-reported health to classify
individuals into good and bad health and estimate the probabilities of
transitioning between the two.  Specifically, we use a logistic regression
to estimate the probability of being in bad health in the next period:
\begin{equation}
  \pi(h_{t+1} \mid x_t)
  \;=\; \Lambda_h\!\bigl(Z_h'\,\phi_h\bigr),
  \label{eq:health_transition}
\end{equation}
where $Z_h$ includes the current health state and age.  We estimate the
parameters for each education type separately.  We use the estimated
parameters to predict transition rates and simulate with them from the
initial share of healthy individuals.

Table~\ref{tab:health_transition} documents the estimated parameters.
Figures~\ref{fig:health_good} and~\ref{fig:health_bad} show the predicted
and observed shares of women in good and bad health, respectively.

\begin{table}[htbp]
\centering
\caption{Health Transition Parameters}
\label{tab:health_transition}
\input{CareerCosts/tables/publication/stochastic_processes/health_transition_women}
\end{table}

\begin{figure}[htbp]
\centering
\includegraphics[width=0.85\textwidth]{CareerCosts/figures/stochastic_processes/health_prob_healthy_women}
\caption{Share of Women in Good Health}
\label{fig:health_good}
\figurenotes{Observed (dashed) and estimated (solid) share of women in good
health by age and education.  SOEP-Core.}
\end{figure}

\begin{figure}[htbp]
\centering
\includegraphics[width=0.85\textwidth]{CareerCosts/figures/stochastic_processes/health_prob_bad_women}
\caption{Share of Women in Bad Health}
\label{fig:health_bad}
\figurenotes{Observed (dashed) and estimated (solid) share of women in bad
health by age and education.  SOEP-Core.}
\end{figure}

Death is modeled jointly within the health Markov process $h_t$.  In the
event of death, the agent bequeaths all remaining wealth and receives
bequest utility.  The probability of dying depends on the agent's health
state and education type.  We estimate a logistic specification that scales
federal life-table death probabilities by health--education interaction
terms:
\begin{equation}
  p_{\text{death}}(t, h_t, \tau)
  \;=\; p^{\text{LT}}_{\text{death}}(t)\;\cdot\;\exp\!\bigl(\beta_{h_t,\tau}\bigr),
  \label{eq:mortality}
\end{equation}
where $p^{\text{LT}}_{\text{death}}(t)$ is the baseline death probability
from the Federal Statistical Office life table at age~$t$ and
$\beta_{h_t,\tau}$ is the estimated coefficient for health state $h_t$ and
education type~$\tau$.

Table~\ref{tab:mortality} reports the logit estimates and
Figure~\ref{fig:survival} displays the predicted survival probabilities by
health and education.

\begin{table}[htbp]
\centering
\caption{Mortality --- Logit Parameters}
\label{tab:mortality}
\input{CareerCosts/tables/publication/stochastic_processes/mortality_women}
\end{table}

\begin{figure}[htbp]
\centering
\includegraphics[width=0.85\textwidth]{CareerCosts/figures/stochastic_processes/survival_probabilities_women}
\caption{Share of Women Alive}
\label{fig:survival}
\figurenotes{Estimated survival probabilities by health state and education
type.  Good health (solid) and bad health (dotted).  Life-table baseline
adjusted via logit.  SOEP-Core.}
\end{figure}


% ===================================================================
\section{Parental Health and Care Demand}
\label{sec:parental_health_care}
% ===================================================================

The agent's caregiving decisions depend on her mother's health and care
needs.  We model three interconnected processes: the mother's mortality
(Section~\ref{sec:mother_mortality}), the mother's functional limitations
measured by ADL categories (Section~\ref{sec:mother_adl}), and the
resulting care demand (Section~\ref{sec:care_demand}).  The caregiving type
(A or B) determines whether the agent is a potential informal caregiver;
type-A agents cannot provide informal care themselves, while type-B agents
face the choice between informal and formal care arrangements when care is
demanded.

Additionally, we estimate the probability that another family member
provides informal care when the agent's caregiving type is~A
(Section~\ref{sec:exog_care_supply}).


% -------------------------------------------------------------------
\subsection{Mother's Mortality}
\label{sec:mother_mortality}
% -------------------------------------------------------------------

The mother's survival is modeled via age-specific death probabilities drawn
from federal life tables published by the Federal Statistical Office.  The
mother's age is inferred from the agent's age and the mean
mother--daughter age difference, which we compute by education type from
the SOEP structural sample:
\begin{equation}
  \text{age}^{\text{mother}}_{t}
  \;=\; t + \overline{\Delta}_{\tau}^{\text{mother}} + \delta,
  \label{eq:mother_age}
\end{equation}
where $\overline{\Delta}_{\tau}^{\text{mother}}$ is the mean
mother--daughter age difference for education type~$\tau$ and $\delta$ is a
fixed offset.

The mother's alive/dead state $m_t$ follows a three-state process: Alive
($m_t = 0$), Recently Dead ($m_t = 1$), and Dead ($m_t = 2$).  The
Recently Dead state persists for exactly one period and triggers the
inheritance process (see Section~\ref{sec:inheritance_prob}).  Transitions
are:
\begin{equation}
  P(m_{t+1} \mid m_t) =
  \begin{cases}
    \bigl[1 - p^{\text{LT}}_{\text{death}}(\text{age}^{\text{mother}}_t),\;
           p^{\text{LT}}_{\text{death}}(\text{age}^{\text{mother}}_t),\; 0\bigr]
    & \text{if } m_t = 0, \\[4pt]
    [0,\; 0,\; 1] & \text{if } m_t = 1, \\[4pt]
    [0,\; 0,\; 1] & \text{if } m_t = 2.
  \end{cases}
  \label{eq:mother_death_transition}
\end{equation}

Table~\ref{tab:mother_mortality} reports the life-table parameters used
for the mother's mortality.

\begin{table}[htbp]
\centering
\caption{Mother's Mortality --- Life-Table Parameters}
\label{tab:mother_mortality}
\input{CareerCosts/tables/publication/stochastic_processes/mother_mortality_women}
\end{table}


% -------------------------------------------------------------------
\subsection{Mother's ADL States}
\label{sec:mother_adl}
% -------------------------------------------------------------------

The mother's functional limitations are tracked via Activities of Daily
Living (ADL) categories, which we collapse into three states: No~ADL
(category~0), Light (category~1), and Intensive (categories~2 and~3
combined).  We estimate a multinomial logistic regression on the Survey of
Health, Ageing and Retirement in Europe (SHARE) parent--child data, using
the mother's lagged ADL category and a cubic polynomial in the mother's
age:
\begin{equation}
  P\!\bigl(\text{ADL}_{t+1} = c \mid \text{ADL}_{t},\,\text{age}^{\text{mother}}_t\bigr)
  \;=\;
  \frac{\exp\!\bigl(X_t'\,\beta_c\bigr)}
       {\sum_{j=0}^{C}\exp\!\bigl(X_t'\,\beta_j\bigr)},
  \label{eq:adl_multinomial}
\end{equation}
where $X_t = \bigl(1,\;\text{age},\;\text{age}^2,\;\text{age}^3,\;
\mathbb{1}[\text{ADL}_t = 1],\;\mathbb{1}[\text{ADL}_t = 2],\;
\mathbb{1}[\text{ADL}_t = 3]\bigr)'$
and the reference category is No~ADL ($c = 0$, normalised to
$\beta_0 = 0$).  The model is estimated separately by sex; we use the
female coefficients for the mother.

Table~\ref{tab:mother_adl} reports the multinomial logit estimates.

\begin{table}[htbp]
\centering
\caption{Mother's ADL Transition --- Multinomial Logit Parameters}
\label{tab:mother_adl}
\input{CareerCosts/tables/publication/stochastic_processes/mother_adl_transition_women}
\end{table}


% -------------------------------------------------------------------
\subsection{Care Demand}
\label{sec:care_demand}
% -------------------------------------------------------------------

Care demand is derived deterministically from the mother's ADL state, the
mother's alive/dead status, and a caregiving age window.  Care can only be
demanded while the mother is alive and the agent's age falls within the
caregiving window (ages~40--70).  Outside this window, or when the mother is
dead, care demand is set to zero regardless of the mother's ADL state.

The care demand probabilities are:
\begin{align}
  p_{\text{none}} &= p_{\text{ADL}_0}
    + \bigl(1 - \mathbb{1}_{\text{window}}\bigr)
    \cdot \bigl(p_{\text{ADL}_1} + p_{\text{ADL}_2}\bigr),
  \label{eq:care_demand_none} \\[2pt]
  p_{\text{light}} &= p_{\text{ADL}_1}
    \cdot \mathbb{1}_{\text{window}},
  \label{eq:care_demand_light} \\[2pt]
  p_{\text{intensive}} &= p_{\text{ADL}_2}
    \cdot \mathbb{1}_{\text{window}},
  \label{eq:care_demand_intensive}
\end{align}
where $p_{\text{ADL}_c}$ are the ADL transition probabilities from
equation~\eqref{eq:adl_multinomial} and the indicator
$\mathbb{1}_{\text{window}}$ equals one if and only if the mother is alive
($m_t = 0$) and the agent's age~$t$ lies in $[40,\,70)$.

Since care demand is fully determined by the mother's ADL process and the
caregiving window, no separate estimation is required.


% -------------------------------------------------------------------
\subsection{Exogenous Care Supply}
\label{sec:exog_care_supply}
% -------------------------------------------------------------------

When care is demanded and the agent's caregiving type is~A (not a potential
informal caregiver), another family member may step in to provide informal
care.  We estimate the probability of receiving such informal care from
other family members via a logistic regression on SOEP data:
\begin{equation}
  P(\text{other informal care} = 1)
  \;=\; \Lambda\!\bigl(\gamma_0 + \gamma_1\,\text{age}
  + \gamma_2\,\text{age}^2
  + \gamma_3\,\text{has\_sister}
  + \gamma_4\,\text{education}\bigr),
  \label{eq:exog_care_supply}
\end{equation}
where the sample is restricted to women whose parent has a care demand.

Table~\ref{tab:exog_care_supply} reports the logit estimates.

\begin{table}[htbp]
\centering
\caption{Exogenous Care Supply --- Logit Parameters}
\label{tab:exog_care_supply}
\input{CareerCosts/tables/publication/stochastic_processes/exog_care_supply_women}
\end{table}


% ===================================================================
\section{Income Processes}
\label{sec:income_processes}
% ===================================================================


% -------------------------------------------------------------------
\subsection{Pension Calculation}
\label{sec:pension}
% -------------------------------------------------------------------

The formula for calculating pension claims in Germany consists of three
parts.  First, there is the pension point value, for which we use the
population-weighted average of East and West German pension point values in
2010.  Second, the pension points accumulated over the working life.  Third,
the deduction factor if the individual retired early.  We track the latter
two through the experience stock $e_t$.

Each individual receives pension points in the ratio of her yearly income
compared to the overall mean wage of all working individuals.  Let $w_m$ be
the mean wage and $h_t$ be the agent's work hours (part-time or full-time).
The pension points earned in any given year are:
\begin{equation}
  \text{PP earned}_t \;=\; \frac{w_t \cdot h_t}{w_m}.
  \label{eq:pension_points_earned}
\end{equation}
Accumulated pension points are the sum over all working years.  We
approximate the cumulative pension points by using the mean experience per
type at age~30 and assuming the agent worked full-time thereafter.

If an agent retires at the statutory retirement age (SRA), she receives the
full pension.  The monthly gross pension is:
\begin{equation}
  y_t(x_t, 0) \;=\; PP(x_t) \times PPV,
  \label{eq:pension_gross}
\end{equation}
where $PPV$ is the pension point value (the 2010 population-weighted
East-West average).

If the agent retires before the SRA, an early-retirement penalty of 3.6\%
per year of early retirement applies:
\begin{equation}
  y_t(x_t, 0) \;=\; PP(x_t) \times PPV \times
  \bigl(1 - 0.036 \times \Delta_{\text{age}}\bigr),
  \label{eq:pension_early_retirement}
\end{equation}
where $\Delta_{\text{age}}$ is the number of years between the actual
retirement age and the SRA.  We track the deduction through the experience
stock: the reduced pension corresponds to an adjusted experience level
$e_t'$ such that the early-retirement penalty is embedded in the state
without explicitly tracking the retirement age.


% -------------------------------------------------------------------
\subsection{Own Wage Process}
\label{sec:own_wage}
% -------------------------------------------------------------------

In the model, agents invest in their human capital by working full-time or
part-time.  We estimate the returns to experience using two-way
fixed-effects regressions on SOEP-Core panel data.  The estimation sample
consists of women over age~30 who work full-time or part-time throughout the
estimation period 2010--2017.  Since time fixed effects absorb aggregate
income growth and inflation, all monetary quantities are expressed in 2010
euros.

We estimate the following equation for each education type~$\tau$ using
observations of wages and experience for each individual~$i$ and time~$t$:
\begin{equation}
  \ln(w_{it})
  \;=\; \gamma_{0,\tau} + \gamma_{\ln,\tau}\,\ln(e_{it} + 1)
  + \xi_i + \mu_t + \varepsilon_{it},
  \label{eq:wage}
\end{equation}
where $\xi_i$ denotes individual fixed effects, $\mu_t$ time fixed effects,
and $\varepsilon_{it}$ is the idiosyncratic wage shock.  Our estimates of
$\gamma_{0,\tau}$ and $\gamma_{\ln,\tau}$ directly correspond to the
structural parameters used in the model.  We cluster standard errors across
individuals and time and estimate the wage process's variance
$\sigma_w^2$.

Table~\ref{tab:own_wage} reports the estimates and
Figure~\ref{fig:own_wage} shows the fit of the predicted wages.

\begin{table}[htbp]
\centering
\caption{Own Wage --- Panel OLS Parameters}
\label{tab:own_wage}
\input{CareerCosts/tables/publication/stochastic_processes/own_wage_women}
\end{table}

\begin{figure}[htbp]
\centering
\includegraphics[width=0.85\textwidth]{CareerCosts/figures/stochastic_processes/own_wage_women}
\caption{Own Wage Fit}
\label{fig:own_wage}
\figurenotes{Observed mean log hourly wages (dashed) and predicted values
(solid) by age and education.  Panel OLS with entity and time fixed
effects.  SOEP-Core, women, 2010--2017.}
\end{figure}


% -------------------------------------------------------------------
\subsection{Partner Income Process}
\label{sec:partner_income}
% -------------------------------------------------------------------

We approximate the partner's income through the agent's own state
variables.  If the agent is single ($p_t = 0$), there is no partner income.
When the agent has a working-age partner ($p_t = 1$), we approximate the
partner's monthly gross wage by OLS of wages onto the agent's age and age
squared.  Unemployed partners are assigned a wage of zero.  The partner
income is therefore a combination of the wages of employed partners and the
zero wages of unemployed partners.  We estimate the approximation separately
by education type.

The specification is:
\begin{equation}
  y^p_t \;=\; \alpha_0 + \alpha_1\,\text{age}_t
  + \alpha_2\,\text{age}_t^2 + u_t.
  \label{eq:partner_wage}
\end{equation}

For retired partners ($p_t = 2$), we use the mean income from the
working-age period and apply a replacement rate of 48\% to obtain the
partner's pension income:
\begin{equation}
  y^{p,\text{ret}}_t
  \;=\; 0.48 \times \overline{y^p}.
  \label{eq:partner_pension}
\end{equation}

Table~\ref{tab:partner_wage} presents the OLS estimates and
Figure~\ref{fig:partner_wage} shows the fit of the approximation during the
partner's working life.

\begin{table}[htbp]
\centering
\caption{Partner Wage --- OLS Parameters}
\label{tab:partner_wage}
\input{CareerCosts/tables/publication/stochastic_processes/partner_wage_women}
\end{table}

\begin{figure}[htbp]
\centering
\includegraphics[width=0.85\textwidth]{CareerCosts/figures/stochastic_processes/partner_wage_women}
\caption{Partner Wage Fit}
\label{fig:partner_wage}
\figurenotes{Observed mean monthly gross partner wage (dashed) and OLS
prediction (solid) by age and education.  SOEP-Core, male partners of
women.}
\end{figure}


% ===================================================================
\section{Parental Bequests and Formal Care Costs}
\label{sec:bequests_care_costs}
% ===================================================================


% -------------------------------------------------------------------
\subsection{Inheritance Probability}
\label{sec:inheritance_prob}
% -------------------------------------------------------------------

We model the probability of receiving a positive inheritance conditional
on the mother's recent death ($m_t = 1$).  The inheritance probability
depends on the agent's age, education type, and the care arrangement in
the period of death.  We estimate the following logistic specification:
\begin{equation}
  P(\text{inheritance} > 0)
  \;=\; \Lambda\!\bigl(\beta_0
  + \beta_{\text{age}}\,\text{age}
  + \beta_{\text{age}^2}\,\text{age}^2
  + \beta_{\text{care}}
  + \beta_{\text{edu}}\,\text{education}\bigr),
  \label{eq:inheritance_prob}
\end{equation}
where $\beta_{\text{care}}$ captures the effect of the care arrangement
(no care, formal care, or informal care) on the probability of receiving a
bequest.

Table~\ref{tab:inheritance_prob} reports the logit estimates.

\begin{table}[htbp]
\centering
\caption{Inheritance Probability --- Logit Parameters}
\label{tab:inheritance_prob}
\input{CareerCosts/tables/publication/stochastic_processes/inheritance_probability_women}
\end{table}


% -------------------------------------------------------------------
\subsection{Inheritance Amount}
\label{sec:inheritance_amount}
% -------------------------------------------------------------------

Conditional on receiving a positive inheritance, we model the amount via a
log-linear OLS regression:
\begin{equation}
  \ln(\text{amount})
  \;=\; \beta_0
  + \beta_{\text{age}}\,\text{age}
  + \beta_{\text{age}^2}\,\text{age}^2
  + \beta_{\text{care}}
  + \beta_{\text{edu}}\,\text{education}
  + \varepsilon,
  \label{eq:inheritance_amount}
\end{equation}
where the care-type coefficients distinguish between light informal care,
intensive informal care, and formal care (no care is the reference).

Table~\ref{tab:inheritance_amount} reports the estimates.

\begin{table}[htbp]
\centering
\caption{Inheritance Amount --- OLS Parameters}
\label{tab:inheritance_amount}
\input{CareerCosts/tables/publication/stochastic_processes/inheritance_amount_women}
\end{table}


% -------------------------------------------------------------------
\subsection{Formal Care Costs}
\label{sec:formal_care_costs}
% -------------------------------------------------------------------

When the agent arranges formal care for her parent, she incurs a financial
cost representing her contribution to the parent's care arrangement.  We
estimate the agent's monthly out-of-pocket contribution to formal care
costs using the SOEP Innovation Sample (SOEP-IS), which contains detailed
information on care arrangements and associated expenditures.

We estimate the following OLS specification, pooled across education
types to avoid selection bias---since selection into formal care is likely
correlated with education, conditioning on education would bias the
estimated costs:
\begin{equation}
  \text{costs}_t
  \;=\; \alpha_0 + \alpha_1\,\text{age}_t
  + \alpha_2\,\text{age}_t^2 + \varepsilon_t.
  \label{eq:formal_care_costs}
\end{equation}

Table~\ref{tab:formal_care_costs} reports the OLS estimates and
Figure~\ref{fig:formal_care_costs} shows the predicted costs by age.

\begin{table}[htbp]
\centering
\caption{Formal Care Costs --- OLS Parameters}
\label{tab:formal_care_costs}
\input{CareerCosts/tables/publication/stochastic_processes/formal_care_costs}
\end{table}

\begin{figure}[htbp]
\centering
\includegraphics[width=0.85\textwidth]{CareerCosts/figures/stochastic_processes/formal_care_costs_women}
\caption{Formal Care Costs by Age}
\label{fig:formal_care_costs}
\figurenotes{Observed mean monthly contribution to formal care costs
(dashed) and OLS prediction (solid) by age.  SOEP-IS.  Pooled across
education levels.}
\end{figure}

%
%  Assumes booktabs, amsmath, amssymb, graphicx are loaded.
%  Table paths: CareerCosts/tables/publication/stochastic_processes/
%  Figure paths: CareerCosts/figures/stochastic_processes/
% ===========================================================================

% Helper command for figure notes
\providecommand{\figurenotes}[1]{%
  \par\vspace{0.4em}%
  \begin{minipage}{\linewidth}%
    \footnotesize\textit{Notes:}\ #1%
  \end{minipage}%
}


% ===================================================================
\section{Variable Overview}
\label{sec:variable_overview}
% ===================================================================

Table~\ref{tab:variable_overview} summarises the state variables, decision
variables, and continuous states of the life-cycle model.  Agents are women
who enter the model at age~30 and can live up to age~100.  Each period
corresponds to one calendar year.

\begin{table}[htbp]
\centering
\caption{Variable Overview}
\label{tab:variable_overview}
\small
\begin{tabular}{lll}
\toprule
\textbf{Variable Name} & \textbf{Symbol} & \textbf{Possible Values} \\
\midrule
\multicolumn{3}{l}{\textbf{Decisions}} \\
\addlinespace
Labor Supply        & $d_t$           & \{0: Retired, 1: Unemployed, \\
                    &                 & \phantom{\{} 2: Part-Time, 3: Full-Time\} \\
Care Arrangement    & $c_t$           & \{0: No Care, 1: Light Informal Care, \\
                    &                 & \phantom{\{} 2: Intensive Informal Care, \\
                    &                 & \phantom{\{} 3: Combination Care, 4: Formal Care\} \\
Consumption         & $\tilde{c}_t$   & $[0,\, a_t]$ \\
\addlinespace
\midrule
\multicolumn{3}{l}{\textbf{Discrete States}} \\
\addlinespace
Age                 & $t$             & $30,\ldots,100$ \\
Education Type      & $\tau$          & \{Low, High\} \\
Caregiving Type     &                 & \{A: Not Potential Caregiver, \\
                    &                 & \phantom{\{} B: Potential Caregiver\} \\
Already Retired     &                 & \{0, 1\} \\
Partner State       & $p_t$           & \{0: Single, 1: Working-Age Partner, \\
                    &                 & \phantom{\{} 2: Retired Partner\} \\
Health State        & $h_t$           & \{0: Good Health, 1: Bad Health, 2: Dead\} \\
Job Offer           & $o_t$           & \{0: No Offer, 1: Job Offer\} \\
Mother Alive        & $m_t$           & \{0: Alive, 1: Recently Dead, 2: Dead\} \\
Mother ADL          &                 & \{0: No ADL, 1: Light, 2: Intensive\} \\
Care Demand         &                 & \{0: None, 1: Light, 2: Intensive\} \\
\addlinespace
\midrule
\multicolumn{3}{l}{\textbf{Continuous States}} \\
\addlinespace
Assets              & $a_t$           & $\mathbb{R}_{\geq 0}$ \\
Work Experience     & $e_t$           & Projection to interval $[0,\,1]$ \\
Taste Shock         & $\epsilon_t(d_t)$ & GEV i.i.d.\ taste shocks \\
\bottomrule
\end{tabular}

\vspace{0.5em}
\begin{minipage}{\linewidth}
\footnotesize\textit{Notes:}
Key derived variables include pension points $PP(x_t)$, the consumption
equivalence scale $n_t(x_t)$, number of children in the household,
hourly wage $w_t(x_t)$, partner income $y^p_t(x_t)$, total household
income $Y_t(x_t, d_t)$, government benefits $B(x_t, d_t)$, and taxes
$T(x_t, d_t)$---all of which are deterministic functions of the state
variables listed above.
\end{minipage}
\end{table}


% ===================================================================
\section{Auxiliary Markov Processes}
\label{sec:auxiliary_markov}
% ===================================================================

All transition processes in this section are estimated on SOEP-Core panel
data, restricted to the structural estimation sample (2010--2017), unless
noted otherwise.  Estimations are performed separately by education type.


% -------------------------------------------------------------------
\subsection{Partner Transitions}
\label{sec:partner_transitions}
% -------------------------------------------------------------------

The partner state $p_t$ evolves stochastically with transition probabilities
that depend on education, age, and the current partner state---none of which
the agent controls directly.  We impose several restrictions that are
informed by the data and simplify the estimation.  First, agents cannot
transition directly from being single to having a retired partner.  Second,
the probability of becoming single is independent of whether the partner is
currently of working age or retired; that is, separation rates are the same
for both partnered states.  Third, there is no transition from a retired
partner back to a working-age partner, reflecting the absorbing nature of
retirement.  Fourth, all partners are assumed to be retired by age~75;
beyond that age, single agents remain single and partnered agents remain
partnered.  This last assumption proxies for the empirical observation that
most late-life partner-state transitions are attributable to spousal death,
which we capture via the widow pension.

Formally, the transition of the partner state is given by:
\begin{equation}
  \pi(p_{t+1} \mid x_t) \;=\; \Lambda_p\!\bigl(Z_p(x_t)'\,\phi_p\bigr),
  \label{eq:partner_transition}
\end{equation}
where $\Lambda_p$ is the multinomial logistic distribution function,
providing transition probabilities over the three partner states.  The
characteristics in $Z_p(x_t)$ are education, age, and the current partner
state.  We follow a two-step estimation strategy: first, we estimate
marriage and separation probabilities, and second, we estimate the
transition of the partner into retirement.  We estimate partner transitions
separately for each of the two education types.

Table~\ref{tab:partner_transition} reports the estimated transition
probabilities.  Figures~\ref{fig:partner_state_no_partner}--\ref{fig:partner_state_retired}
illustrate the model fit by comparing simulated shares with the empirical
counterparts.

\begin{table}[htbp]
\centering
\caption{Partner Transition Probabilities}
\label{tab:partner_transition}
\input{CareerCosts/tables/publication/stochastic_processes/partner_transition_women}
\end{table}

\begin{figure}[htbp]
\centering
\includegraphics[width=0.85\textwidth]{CareerCosts/figures/stochastic_processes/partner_state_no_partner_women}
\caption{Share of Women Without a Partner}
\label{fig:partner_state_no_partner}
\figurenotes{Observed shares (dashed) and estimated shares (solid) by age
and education type.  Estimation on SOEP-Core, women only.}
\end{figure}

\begin{figure}[htbp]
\centering
\includegraphics[width=0.85\textwidth]{CareerCosts/figures/stochastic_processes/partner_state_working_partner_women}
\caption{Share of Women With a Working-Age Partner}
\label{fig:partner_state_working}
\figurenotes{Observed shares (dashed) and estimated shares (solid) by age
and education type.  Estimation on SOEP-Core, women only.}
\end{figure}

\begin{figure}[htbp]
\centering
\includegraphics[width=0.85\textwidth]{CareerCosts/figures/stochastic_processes/partner_state_retired_partner_women}
\caption{Share of Women With a Retired Partner}
\label{fig:partner_state_retired}
\figurenotes{Observed shares (dashed) and estimated shares (solid) by age
and education type.  Estimation on SOEP-Core, women only.}
\end{figure}


% -------------------------------------------------------------------
\subsection{Number of Children in Household}
\label{sec:children}
% -------------------------------------------------------------------

The agent's partner state, together with age and education type, determines
the number of children living in the household.  The number of children
directly enters the construction of the consumption equivalence scale and,
conditional on the agent being employed, affects utility through the
preference parameters for children.  We approximate the number of children
in the household via OLS, estimating the following specification separately
by education type and partner status:
\begin{equation}
  n_{t}^{\text{child}} \;=\; \alpha_{0} + \alpha_{1}\,\text{period}_{t}
  + \alpha_{2}\,\text{period}_{t}^{2} + u_{t},
  \label{eq:children}
\end{equation}
where $\text{period}_{t} = t - 30$ denotes model age.

Table~\ref{tab:children} presents the OLS estimates.
Figures~\ref{fig:children_single} and~\ref{fig:children_partnered}
show the estimation fit for single and partnered women, respectively.

\begin{table}[htbp]
\centering
\caption{Number of Children in Household --- OLS Parameters}
\label{tab:children}
\input{CareerCosts/tables/publication/stochastic_processes/number_of_children_women}
\end{table}

\begin{figure}[htbp]
\centering
\includegraphics[width=0.85\textwidth]{CareerCosts/figures/stochastic_processes/children_single_women}
\caption{Number of Children --- Single Women}
\label{fig:children_single}
\figurenotes{Observed (dashed) and estimated (solid) mean number of children
in the household by age and education.  SOEP-Core, single women.}
\end{figure}

\begin{figure}[htbp]
\centering
\includegraphics[width=0.85\textwidth]{CareerCosts/figures/stochastic_processes/children_partnered_women}
\caption{Number of Children --- Partnered Women}
\label{fig:children_partnered}
\figurenotes{Observed (dashed) and estimated (solid) mean number of children
in the household by age and education.  SOEP-Core, partnered women.}
\end{figure}


% -------------------------------------------------------------------
\subsection{Job Offer and Separation}
\label{sec:job_offer_separation}
% -------------------------------------------------------------------

The job-offer state $o_t$ governs the agent's ability to choose
employment; the agent can work part-time or full-time only if $o_t = 1$.
If the job-offer state equals zero and the agent has not yet retired, she
is forced into unemployment in the current period.  We incorporate two
independent processes into the job-offer state: job destruction for those
currently employed and job offers for those currently non-employed.

Job destruction probabilities can be directly identified from the SOEP
using the question on the reason for job departure.  We classify
individuals in job-offer state~0 who did not voluntarily leave their
jobs as having experienced an involuntary separation.  The transition
probability for job separation conditional on current employment is given
by:
\begin{equation}
  \pi(o_{t+1} = 0 \mid x_t,\, d_t \in \{2,3\})
  \;=\; \Lambda_{\text{sep}}\!\bigl(Z_{\text{sep}}'\,\phi_{\text{sep}}\bigr),
  \label{eq:job_separation}
\end{equation}
where $\Lambda_{\text{sep}}$ is the logistic distribution function and
$Z_{\text{sep}}$ includes education, age, and a constant.  We separately
estimate separation probabilities by education type.  We restrict the
estimation sample to ages~30--65 to ensure sufficient observational power
and assume that job separation rates remain constant after age~65 until the
mandatory retirement age.

Table~\ref{tab:job_separation} reports the logit estimates and
Figure~\ref{fig:job_separation} shows the predicted separation
probabilities against the observed rates.

\begin{table}[htbp]
\centering
\caption{Job Separation --- Logit Parameters}
\label{tab:job_separation}
\input{CareerCosts/tables/publication/stochastic_processes/job_separation_women}
\end{table}

\begin{figure}[htbp]
\centering
\includegraphics[width=0.85\textwidth]{CareerCosts/figures/stochastic_processes/job_separation_women}
\caption{Job Separation Probabilities}
\label{fig:job_separation}
\figurenotes{Observed (dashed) and estimated (solid) job separation
probabilities by age and education.  Logit estimation using the
\textit{plb0304\_h} variable from the pl sample in the SOEP.}
\end{figure}

The job-offer process for non-employed agents is estimated internally via
the Method of Simulated Moments (MSM).  If the agent chooses non-employment
during a given period, the job-offer process determines the probability of
having an employment choice available in the next period ($o_{t+1} = 1$).


% -------------------------------------------------------------------
\subsection{Health and Death}
\label{sec:health_death}
% -------------------------------------------------------------------

The health process $h_t$ tracks two living states---Good Health and Bad
Health---together with Death.  Survival probabilities are type-specific and
depend on health and age.  This motivates the use of a joint Markov process
to model the stochastic transitions of health and survival simultaneously.

We use the SOEP-Core question on self-reported health to classify
individuals into good and bad health and estimate the probabilities of
transitioning between the two.  Specifically, we use a logistic regression
to estimate the probability of being in bad health in the next period:
\begin{equation}
  \pi(h_{t+1} \mid x_t)
  \;=\; \Lambda_h\!\bigl(Z_h'\,\phi_h\bigr),
  \label{eq:health_transition}
\end{equation}
where $Z_h$ includes the current health state and age.  We estimate the
parameters for each education type separately.  We use the estimated
parameters to predict transition rates and simulate with them from the
initial share of healthy individuals.

Table~\ref{tab:health_transition} documents the estimated parameters.
Figures~\ref{fig:health_good} and~\ref{fig:health_bad} show the predicted
and observed shares of women in good and bad health, respectively.

\begin{table}[htbp]
\centering
\caption{Health Transition Parameters}
\label{tab:health_transition}
\input{CareerCosts/tables/publication/stochastic_processes/health_transition_women}
\end{table}

\begin{figure}[htbp]
\centering
\includegraphics[width=0.85\textwidth]{CareerCosts/figures/stochastic_processes/health_prob_healthy_women}
\caption{Share of Women in Good Health}
\label{fig:health_good}
\figurenotes{Observed (dashed) and estimated (solid) share of women in good
health by age and education.  SOEP-Core.}
\end{figure}

\begin{figure}[htbp]
\centering
\includegraphics[width=0.85\textwidth]{CareerCosts/figures/stochastic_processes/health_prob_bad_women}
\caption{Share of Women in Bad Health}
\label{fig:health_bad}
\figurenotes{Observed (dashed) and estimated (solid) share of women in bad
health by age and education.  SOEP-Core.}
\end{figure}

Death is modeled jointly within the health Markov process $h_t$.  In the
event of death, the agent bequeaths all remaining wealth and receives
bequest utility.  The probability of dying depends on the agent's health
state and education type.  We estimate a logistic specification that scales
federal life-table death probabilities by health--education interaction
terms:
\begin{equation}
  p_{\text{death}}(t, h_t, \tau)
  \;=\; p^{\text{LT}}_{\text{death}}(t)\;\cdot\;\exp\!\bigl(\beta_{h_t,\tau}\bigr),
  \label{eq:mortality}
\end{equation}
where $p^{\text{LT}}_{\text{death}}(t)$ is the baseline death probability
from the Federal Statistical Office life table at age~$t$ and
$\beta_{h_t,\tau}$ is the estimated coefficient for health state $h_t$ and
education type~$\tau$.

Table~\ref{tab:mortality} reports the logit estimates and
Figure~\ref{fig:survival} displays the predicted survival probabilities by
health and education.

\begin{table}[htbp]
\centering
\caption{Mortality --- Logit Parameters}
\label{tab:mortality}
\input{CareerCosts/tables/publication/stochastic_processes/mortality_women}
\end{table}

\begin{figure}[htbp]
\centering
\includegraphics[width=0.85\textwidth]{CareerCosts/figures/stochastic_processes/survival_probabilities_women}
\caption{Share of Women Alive}
\label{fig:survival}
\figurenotes{Estimated survival probabilities by health state and education
type.  Good health (solid) and bad health (dotted).  Life-table baseline
adjusted via logit.  SOEP-Core.}
\end{figure}


% ===================================================================
\section{Parental Health and Care Demand}
\label{sec:parental_health_care}
% ===================================================================

The agent's caregiving decisions depend on her mother's health and care
needs.  We model three interconnected processes: the mother's mortality
(Section~\ref{sec:mother_mortality}), the mother's functional limitations
measured by ADL categories (Section~\ref{sec:mother_adl}), and the
resulting care demand (Section~\ref{sec:care_demand}).  The caregiving type
(A or B) determines whether the agent is a potential informal caregiver;
type-A agents cannot provide informal care themselves, while type-B agents
face the choice between informal and formal care arrangements when care is
demanded.

Additionally, we estimate the probability that another family member
provides informal care when the agent's caregiving type is~A
(Section~\ref{sec:exog_care_supply}).


% -------------------------------------------------------------------
\subsection{Mother's Mortality}
\label{sec:mother_mortality}
% -------------------------------------------------------------------

The mother's survival is modeled via age-specific death probabilities drawn
from federal life tables published by the Federal Statistical Office.  The
mother's age is inferred from the agent's age and the mean
mother--daughter age difference, which we compute by education type from
the SOEP structural sample:
\begin{equation}
  \text{age}^{\text{mother}}_{t}
  \;=\; t + \overline{\Delta}_{\tau}^{\text{mother}} + \delta,
  \label{eq:mother_age}
\end{equation}
where $\overline{\Delta}_{\tau}^{\text{mother}}$ is the mean
mother--daughter age difference for education type~$\tau$ and $\delta$ is a
fixed offset.

The mother's alive/dead state $m_t$ follows a three-state process: Alive
($m_t = 0$), Recently Dead ($m_t = 1$), and Dead ($m_t = 2$).  The
Recently Dead state persists for exactly one period and triggers the
inheritance process (see Section~\ref{sec:inheritance_prob}).  Transitions
are:
\begin{equation}
  P(m_{t+1} \mid m_t) =
  \begin{cases}
    \bigl[1 - p^{\text{LT}}_{\text{death}}(\text{age}^{\text{mother}}_t),\;
           p^{\text{LT}}_{\text{death}}(\text{age}^{\text{mother}}_t),\; 0\bigr]
    & \text{if } m_t = 0, \\[4pt]
    [0,\; 0,\; 1] & \text{if } m_t = 1, \\[4pt]
    [0,\; 0,\; 1] & \text{if } m_t = 2.
  \end{cases}
  \label{eq:mother_death_transition}
\end{equation}

Table~\ref{tab:mother_mortality} reports the life-table parameters used
for the mother's mortality.

\begin{table}[htbp]
\centering
\caption{Mother's Mortality --- Life-Table Parameters}
\label{tab:mother_mortality}
\input{CareerCosts/tables/publication/stochastic_processes/mother_mortality_women}
\end{table}


% -------------------------------------------------------------------
\subsection{Mother's ADL States}
\label{sec:mother_adl}
% -------------------------------------------------------------------

The mother's functional limitations are tracked via Activities of Daily
Living (ADL) categories, which we collapse into three states: No~ADL
(category~0), Light (category~1), and Intensive (categories~2 and~3
combined).  We estimate a multinomial logistic regression on the Survey of
Health, Ageing and Retirement in Europe (SHARE) parent--child data, using
the mother's lagged ADL category and a cubic polynomial in the mother's
age:
\begin{equation}
  P\!\bigl(\text{ADL}_{t+1} = c \mid \text{ADL}_{t},\,\text{age}^{\text{mother}}_t\bigr)
  \;=\;
  \frac{\exp\!\bigl(X_t'\,\beta_c\bigr)}
       {\sum_{j=0}^{C}\exp\!\bigl(X_t'\,\beta_j\bigr)},
  \label{eq:adl_multinomial}
\end{equation}
where $X_t = \bigl(1,\;\text{age},\;\text{age}^2,\;\text{age}^3,\;
\mathbb{1}[\text{ADL}_t = 1],\;\mathbb{1}[\text{ADL}_t = 2],\;
\mathbb{1}[\text{ADL}_t = 3]\bigr)'$
and the reference category is No~ADL ($c = 0$, normalised to
$\beta_0 = 0$).  The model is estimated separately by sex; we use the
female coefficients for the mother.

Table~\ref{tab:mother_adl} reports the multinomial logit estimates.

\begin{table}[htbp]
\centering
\caption{Mother's ADL Transition --- Multinomial Logit Parameters}
\label{tab:mother_adl}
\input{CareerCosts/tables/publication/stochastic_processes/mother_adl_transition_women}
\end{table}


% -------------------------------------------------------------------
\subsection{Care Demand}
\label{sec:care_demand}
% -------------------------------------------------------------------

Care demand is derived deterministically from the mother's ADL state, the
mother's alive/dead status, and a caregiving age window.  Care can only be
demanded while the mother is alive and the agent's age falls within the
caregiving window (ages~40--70).  Outside this window, or when the mother is
dead, care demand is set to zero regardless of the mother's ADL state.

The care demand probabilities are:
\begin{align}
  p_{\text{none}} &= p_{\text{ADL}_0}
    + \bigl(1 - \mathbb{1}_{\text{window}}\bigr)
    \cdot \bigl(p_{\text{ADL}_1} + p_{\text{ADL}_2}\bigr),
  \label{eq:care_demand_none} \\[2pt]
  p_{\text{light}} &= p_{\text{ADL}_1}
    \cdot \mathbb{1}_{\text{window}},
  \label{eq:care_demand_light} \\[2pt]
  p_{\text{intensive}} &= p_{\text{ADL}_2}
    \cdot \mathbb{1}_{\text{window}},
  \label{eq:care_demand_intensive}
\end{align}
where $p_{\text{ADL}_c}$ are the ADL transition probabilities from
equation~\eqref{eq:adl_multinomial} and the indicator
$\mathbb{1}_{\text{window}}$ equals one if and only if the mother is alive
($m_t = 0$) and the agent's age~$t$ lies in $[40,\,70)$.

Since care demand is fully determined by the mother's ADL process and the
caregiving window, no separate estimation is required.


% -------------------------------------------------------------------
\subsection{Exogenous Care Supply}
\label{sec:exog_care_supply}
% -------------------------------------------------------------------

When care is demanded and the agent's caregiving type is~A (not a potential
informal caregiver), another family member may step in to provide informal
care.  We estimate the probability of receiving such informal care from
other family members via a logistic regression on SOEP data:
\begin{equation}
  P(\text{other informal care} = 1)
  \;=\; \Lambda\!\bigl(\gamma_0 + \gamma_1\,\text{age}
  + \gamma_2\,\text{age}^2
  + \gamma_3\,\text{has\_sister}
  + \gamma_4\,\text{education}\bigr),
  \label{eq:exog_care_supply}
\end{equation}
where the sample is restricted to women whose parent has a care demand.

Table~\ref{tab:exog_care_supply} reports the logit estimates.

\begin{table}[htbp]
\centering
\caption{Exogenous Care Supply --- Logit Parameters}
\label{tab:exog_care_supply}
\input{CareerCosts/tables/publication/stochastic_processes/exog_care_supply_women}
\end{table}


% ===================================================================
\section{Income Processes}
\label{sec:income_processes}
% ===================================================================


% -------------------------------------------------------------------
\subsection{Pension Calculation}
\label{sec:pension}
% -------------------------------------------------------------------

The formula for calculating pension claims in Germany consists of three
parts.  First, there is the pension point value, for which we use the
population-weighted average of East and West German pension point values in
2010.  Second, the pension points accumulated over the working life.  Third,
the deduction factor if the individual retired early.  We track the latter
two through the experience stock $e_t$.

Each individual receives pension points in the ratio of her yearly income
compared to the overall mean wage of all working individuals.  Let $w_m$ be
the mean wage and $h_t$ be the agent's work hours (part-time or full-time).
The pension points earned in any given year are:
\begin{equation}
  \text{PP earned}_t \;=\; \frac{w_t \cdot h_t}{w_m}.
  \label{eq:pension_points_earned}
\end{equation}
Accumulated pension points are the sum over all working years.  We
approximate the cumulative pension points by using the mean experience per
type at age~30 and assuming the agent worked full-time thereafter.

If an agent retires at the statutory retirement age (SRA), she receives the
full pension.  The monthly gross pension is:
\begin{equation}
  y_t(x_t, 0) \;=\; PP(x_t) \times PPV,
  \label{eq:pension_gross}
\end{equation}
where $PPV$ is the pension point value (the 2010 population-weighted
East-West average).

If the agent retires before the SRA, an early-retirement penalty of 3.6\%
per year of early retirement applies:
\begin{equation}
  y_t(x_t, 0) \;=\; PP(x_t) \times PPV \times
  \bigl(1 - 0.036 \times \Delta_{\text{age}}\bigr),
  \label{eq:pension_early_retirement}
\end{equation}
where $\Delta_{\text{age}}$ is the number of years between the actual
retirement age and the SRA.  We track the deduction through the experience
stock: the reduced pension corresponds to an adjusted experience level
$e_t'$ such that the early-retirement penalty is embedded in the state
without explicitly tracking the retirement age.


% -------------------------------------------------------------------
\subsection{Own Wage Process}
\label{sec:own_wage}
% -------------------------------------------------------------------

In the model, agents invest in their human capital by working full-time or
part-time.  We estimate the returns to experience using two-way
fixed-effects regressions on SOEP-Core panel data.  The estimation sample
consists of women over age~30 who work full-time or part-time throughout the
estimation period 2010--2017.  Since time fixed effects absorb aggregate
income growth and inflation, all monetary quantities are expressed in 2010
euros.

We estimate the following equation for each education type~$\tau$ using
observations of wages and experience for each individual~$i$ and time~$t$:
\begin{equation}
  \ln(w_{it})
  \;=\; \gamma_{0,\tau} + \gamma_{\ln,\tau}\,\ln(e_{it} + 1)
  + \xi_i + \mu_t + \varepsilon_{it},
  \label{eq:wage}
\end{equation}
where $\xi_i$ denotes individual fixed effects, $\mu_t$ time fixed effects,
and $\varepsilon_{it}$ is the idiosyncratic wage shock.  Our estimates of
$\gamma_{0,\tau}$ and $\gamma_{\ln,\tau}$ directly correspond to the
structural parameters used in the model.  We cluster standard errors across
individuals and time and estimate the wage process's variance
$\sigma_w^2$.

Table~\ref{tab:own_wage} reports the estimates and
Figure~\ref{fig:own_wage} shows the fit of the predicted wages.

\begin{table}[htbp]
\centering
\caption{Own Wage --- Panel OLS Parameters}
\label{tab:own_wage}
\input{CareerCosts/tables/publication/stochastic_processes/own_wage_women}
\end{table}

\begin{figure}[htbp]
\centering
\includegraphics[width=0.85\textwidth]{CareerCosts/figures/stochastic_processes/own_wage_women}
\caption{Own Wage Fit}
\label{fig:own_wage}
\figurenotes{Observed mean log hourly wages (dashed) and predicted values
(solid) by age and education.  Panel OLS with entity and time fixed
effects.  SOEP-Core, women, 2010--2017.}
\end{figure}


% -------------------------------------------------------------------
\subsection{Partner Income Process}
\label{sec:partner_income}
% -------------------------------------------------------------------

We approximate the partner's income through the agent's own state
variables.  If the agent is single ($p_t = 0$), there is no partner income.
When the agent has a working-age partner ($p_t = 1$), we approximate the
partner's monthly gross wage by OLS of wages onto the agent's age and age
squared.  Unemployed partners are assigned a wage of zero.  The partner
income is therefore a combination of the wages of employed partners and the
zero wages of unemployed partners.  We estimate the approximation separately
by education type.

The specification is:
\begin{equation}
  y^p_t \;=\; \alpha_0 + \alpha_1\,\text{age}_t
  + \alpha_2\,\text{age}_t^2 + u_t.
  \label{eq:partner_wage}
\end{equation}

For retired partners ($p_t = 2$), we use the mean income from the
working-age period and apply a replacement rate of 48\% to obtain the
partner's pension income:
\begin{equation}
  y^{p,\text{ret}}_t
  \;=\; 0.48 \times \overline{y^p}.
  \label{eq:partner_pension}
\end{equation}

Table~\ref{tab:partner_wage} presents the OLS estimates and
Figure~\ref{fig:partner_wage} shows the fit of the approximation during the
partner's working life.

\begin{table}[htbp]
\centering
\caption{Partner Wage --- OLS Parameters}
\label{tab:partner_wage}
\input{CareerCosts/tables/publication/stochastic_processes/partner_wage_women}
\end{table}

\begin{figure}[htbp]
\centering
\includegraphics[width=0.85\textwidth]{CareerCosts/figures/stochastic_processes/partner_wage_women}
\caption{Partner Wage Fit}
\label{fig:partner_wage}
\figurenotes{Observed mean monthly gross partner wage (dashed) and OLS
prediction (solid) by age and education.  SOEP-Core, male partners of
women.}
\end{figure}


% ===================================================================
\section{Parental Bequests and Formal Care Costs}
\label{sec:bequests_care_costs}
% ===================================================================


% -------------------------------------------------------------------
\subsection{Inheritance Probability}
\label{sec:inheritance_prob}
% -------------------------------------------------------------------

We model the probability of receiving a positive inheritance conditional
on the mother's recent death ($m_t = 1$).  The inheritance probability
depends on the agent's age, education type, and the care arrangement in
the period of death.  We estimate the following logistic specification:
\begin{equation}
  P(\text{inheritance} > 0)
  \;=\; \Lambda\!\bigl(\beta_0
  + \beta_{\text{age}}\,\text{age}
  + \beta_{\text{age}^2}\,\text{age}^2
  + \beta_{\text{care}}
  + \beta_{\text{edu}}\,\text{education}\bigr),
  \label{eq:inheritance_prob}
\end{equation}
where $\beta_{\text{care}}$ captures the effect of the care arrangement
(no care, formal care, or informal care) on the probability of receiving a
bequest.

Table~\ref{tab:inheritance_prob} reports the logit estimates.

\begin{table}[htbp]
\centering
\caption{Inheritance Probability --- Logit Parameters}
\label{tab:inheritance_prob}
\input{CareerCosts/tables/publication/stochastic_processes/inheritance_probability_women}
\end{table}


% -------------------------------------------------------------------
\subsection{Inheritance Amount}
\label{sec:inheritance_amount}
% -------------------------------------------------------------------

Conditional on receiving a positive inheritance, we model the amount via a
log-linear OLS regression:
\begin{equation}
  \ln(\text{amount})
  \;=\; \beta_0
  + \beta_{\text{age}}\,\text{age}
  + \beta_{\text{age}^2}\,\text{age}^2
  + \beta_{\text{care}}
  + \beta_{\text{edu}}\,\text{education}
  + \varepsilon,
  \label{eq:inheritance_amount}
\end{equation}
where the care-type coefficients distinguish between light informal care,
intensive informal care, and formal care (no care is the reference).

Table~\ref{tab:inheritance_amount} reports the estimates.

\begin{table}[htbp]
\centering
\caption{Inheritance Amount --- OLS Parameters}
\label{tab:inheritance_amount}
\input{CareerCosts/tables/publication/stochastic_processes/inheritance_amount_women}
\end{table}


% -------------------------------------------------------------------
\subsection{Formal Care Costs}
\label{sec:formal_care_costs}
% -------------------------------------------------------------------

When the agent arranges formal care for her parent, she incurs a financial
cost representing her contribution to the parent's care arrangement.  We
estimate the agent's monthly out-of-pocket contribution to formal care
costs using the SOEP Innovation Sample (SOEP-IS), which contains detailed
information on care arrangements and associated expenditures.

We estimate the following OLS specification, pooled across education
types to avoid selection bias---since selection into formal care is likely
correlated with education, conditioning on education would bias the
estimated costs:
\begin{equation}
  \text{costs}_t
  \;=\; \alpha_0 + \alpha_1\,\text{age}_t
  + \alpha_2\,\text{age}_t^2 + \varepsilon_t.
  \label{eq:formal_care_costs}
\end{equation}

Table~\ref{tab:formal_care_costs} reports the OLS estimates and
Figure~\ref{fig:formal_care_costs} shows the predicted costs by age.

\begin{table}[htbp]
\centering
\caption{Formal Care Costs --- OLS Parameters}
\label{tab:formal_care_costs}
\input{CareerCosts/tables/publication/stochastic_processes/formal_care_costs}
\end{table}

\begin{figure}[htbp]
\centering
\includegraphics[width=0.85\textwidth]{CareerCosts/figures/stochastic_processes/formal_care_costs_women}
\caption{Formal Care Costs by Age}
\label{fig:formal_care_costs}
\figurenotes{Observed mean monthly contribution to formal care costs
(dashed) and OLS prediction (solid) by age.  SOEP-IS.  Pooled across
education levels.}
\end{figure}
